%%%%%%%%%%%%%%%%%%%%%%%%%%%%%%%%%%%%%%%%%%%%%%%%%%%%%%%%%%%%%%%%%%%%%%%%%%%
%% Reasoning-aware Adaptive Routing for Hybrid GraphRAG
%% Target journal: Discover Artificial Intelligence (Springer Nature, #44163)
%% Article type in Snapp: "Research" (alternatively "Methodology" -- see REVISION_NOTES)
%%
%% SINGLE-FILE (flattened) manuscript, per Springer Nature's one-file rule.
%% English-only version (all Vietnamese surface text translated to English;
%% prompts reproduced in English only).
%%
%% Template: official Springer Nature sn-jnl.cls with sn-mathphys-num.bst.
%% Download from:
%% https://www.springernature.com/gp/authors/campaigns/latex-author-support
%% Place sn-jnl.cls + sn-mathphys-num.bst in this project root.
%%%%%%%%%%%%%%%%%%%%%%%%%%%%%%%%%%%%%%%%%%%%%%%%%%%%%%%%%%%%%%%%%%%%%%%%%%%

\documentclass[pdflatex,sn-mathphys-num]{sn-jnl}

% ---------------------------------------------------------------
% Packages (sn-jnl preloads amsmath/graphicx-like essentials; keep minimal)
% ---------------------------------------------------------------
\usepackage{graphicx}
\usepackage{multirow}
\usepackage{amsmath,amssymb,amsfonts}
\usepackage{mathtools}
\usepackage{booktabs}
\usepackage{tabularx}
\usepackage{array}
\usepackage{ragged2e}
\usepackage{algorithm}
\usepackage{algpseudocode}
\usepackage[T1]{fontenc}   % standard Latin encoding (manuscript is English-only)
\usepackage{lmodern}
\usepackage{textcomp}
\usepackage{url}

% Column types for description tables
\newcolumntype{Y}{>{\RaggedRight\arraybackslash}X}
\newcolumntype{L}[1]{>{\RaggedRight\arraybackslash}p{#1}}

% Helper to typeset cross-validation confidence half-widths (mean +/- std)
\newcommand{\ci}[1]{\ensuremath{\scriptstyle\pm #1}}

% Keep submission diagnostics focused on actionable errors.
\hypersetup{hypertexnames=false,bookmarksdepth=2}
\hbadness=10000
\vbadness=10000
\hfuzz=100pt

% ---------------------------------------------------------------
\begin{document}

\title[Reasoning-aware Adaptive Routing for Hybrid GraphRAG]{Reasoning-aware Adaptive Routing for Hybrid GraphRAG in Vietnamese Legal Question Answering}

\author[1]{\fnm{Cong} \sur{Tran}}\email{congtran@ptit.edu.vn}

\author*[2]{\fnm{Nguyen Dinh} \sur{Thanh}}\email{thanh.nd225670@sis.hust.edu.vn}

\author[2]{\fnm{Pham Van} \sur{Hai}}\email{haipv@soict.hust.edu.vn}

\affil[1]{\orgname{Posts and Telecommunications Institute of Technology (PTIT)}, \city{Hanoi}, \country{Vietnam}}

\affil[2]{\orgdiv{School of Information and Communication Technology}, \orgname{Hanoi University of Science and Technology}, \city{Hanoi}, \country{Vietnam}}

% ---------------------------------------------------------------
% Abstract: MUST be < 250 words (Discover AI requirement).
% ---------------------------------------------------------------
\abstract{Retrieval-augmented generation for legal question answering commonly
commits to a single retrieval strategy, yet legal questions are heterogeneous:
some require only direct textual lookup, others demand reasoning over
structural relations such as amendment and cross-reference, and still others
are ambiguous and should be clarified rather than answered. Since applying one
fixed strategy to all of them is inherently suboptimal, we propose a
reasoning-aware adaptive router for a Hybrid GraphRAG system that, before any
evidence is retrieved, dispatches each query to dense retrieval, graph
traversal, hybrid reasoning, or clarification. Our router performs two stages:
a lightweight classifier resolves clear cases at negligible cost, while a
selective large language model verifier reasons over the uncertain, ambiguous,
and relation-heavy ones; a deterministic history resolver settles
conversational references beforehand. Through experimental evaluation on a
Vietnamese legal corpus, we demonstrate a) the routing accuracy of the
classifier, b) the causal contribution of the verifier as a chain-of-thought
step, c) the safety of clarification under ambiguity, d) the cost--accuracy
trade-off of selective verification, and e) the behaviour under multi-turn
conversation. A live demonstration of the deployed system is available at
\url{https://naturals-suggested-annex-clothing.trycloudflare.com}, and the
dataset and source code are released at
\url{https://REPLACE-WITH-YOUR-DATA-URL}.}

\keywords{Retrieval-Augmented Generation, GraphRAG, Hybrid RAG, Adaptive
Routing, Vietnamese Legal QA, Knowledge Graphs, Multi-hop Reasoning,
Ambiguity Detection}

\maketitle

% ---------------------------------------------------------------
%% ===== introduction ====================================================
\section{Introduction}\label{sec:intro}
% ===============================================================

\subsection{Background and Motivation}\label{subsec:background}

Retrieval-augmented generation (RAG) has become a standard means of grounding
large language models in external evidence, improving factual reliability
without costly retraining~\cite{lewis2020rag,guu2020realm}. In its most widely
adopted form, RAG relies on dense passage retrieval, embedding the query and
the corpus into a shared vector space and returning the most semantically
similar chunks~\cite{karpukhin2020dpr}. This design is effective for factoid
questions whose answer lies in a single passage, and much of the recent
progress in question answering has been built upon it.

In practice and despite its popularity, however, the assumption that a single
dense retriever suffices tends to be overly simplistic, as the information
need behind a query cannot be served by one fixed strategy in many real-world
settings. When the answer depends on relations among legal provisions, on
document structure, or on chains of evidence spanning several documents, dense
retrieval cannot follow those connections, because the relations are absent
from the vector space~\cite{trivedi2023interleaving}. Graph-based retrieval
was introduced precisely to address this gap by organising the corpus as a
knowledge graph and traversing explicit, typed
relationships~\cite{edge2024graphrag,wang2024kgp,gutierrez2024hipporag}; yet
graph traversal is more expensive than dense lookup, and applying it to every
query wastes computation on the many questions that a single article already
answers. Vietnamese legal question answering exposes this tension sharply: one
system must serve both simple article-level definitions and multi-document
questions involving amendment, replacement, and cross-document
reference~\cite{chalkidis2020legalbert,zhong2020nlp_law,sovrano2020legalkg},
while some questions are ambiguous and admit no safe answer at all until the
user is asked to clarify.

When confronted with such heterogeneous queries, a naturally arising question
is therefore how to decide \emph{which} retrieval strategy a given query
requires, and to do so \emph{before} any evidence is retrieved. Recent studies
have addressed a related question by routing among model sizes or retrieval
depths according to query complexity~\cite{jeong2024adaptive,chen2023frugalgpt,ong2024routellm};
however, they select \emph{how much} to retrieve rather than \emph{which
evidence mechanism} to use, and the prior decision of whether a query needs
relational traversal at all is left unaddressed. A separate and
safety-critical gap concerns ambiguity: when a conversational referent cannot
be resolved, forcing retrieval grounds the answer in an arbitrary document and
returns a confidently wrong answer~\cite{min2020ambigqa,elgohary2019canard},
so the option to clarify must itself be part of the routing decision. This
motivates us to design a router that predicts the appropriate evidence action
directly from the query and its conversational context, treating the strategy
as a property of the query rather than a fixed pipeline setting.

\subsection{Main Contributions}\label{subsec:contributions}

In this paper, we propose a reasoning-aware adaptive router for a Hybrid
GraphRAG system that, before retrieval, dispatches each Vietnamese legal query
to one of four evidence actions---dense retrieval, graph traversal, hybrid
reasoning, or clarification. Rather than running all strategies and merging
their outputs, the router decides upstream through two stages: a lightweight
feature-based classifier resolves the majority of clear cases as a real-time
gateway, and a selective large language model (LLM) verifier is invoked only
for uncertain, ambiguous, or relation-heavy queries, following a cost-aware
cascade principle~\cite{chen2023frugalgpt,ong2024routellm}. A deterministic
history resolver precedes both stages, distinguishing resolved, irrelevant,
and conflicting conversational references so that clarification is triggered
only when the legal target is genuinely underspecified.

Our contributions are four-fold, as summarised below.
\begin{itemize}
\item \textbf{Reasoning-aware routing formulation.} We formulate retrieval
routing for Vietnamese legal question answering as a four-class decision made
\emph{before} evidence retrieval, unifying evidence-mechanism selection and
ambiguity handling---including clarification as a first-class routing
action---within a single framework~\cite{jeong2024adaptive,min2020ambigqa}.
\item \textbf{Two-stage router with a verifier that reasons.} We design a
two-stage router that couples a lightweight classifier with a selectively
invoked LLM verifier, and we identify that this verifier serves not merely as
a route corrector but as an explicit chain-of-thought step that improves the
generated answer, a property we isolate through a controlled oracle-routing
ablation~\cite{chen2023frugalgpt}.
\item \textbf{Deterministic conversational history resolution.} We introduce a
history resolver that classifies a conversational referent as resolved,
irrelevant, or conflicting, enabling the router to unblock retrieval when a
referent is available and to force clarification when it is
absent~\cite{elgohary2019canard}.
\item \textbf{Comprehensive empirical validation.} Through experimental
evaluation on a Vietnamese legal corpus, we demonstrate 1)~the routing
accuracy of the classifier, 2)~the causal contribution of the verifier, 3)~the
safety of clarification under ambiguity, and 4)~the cost--accuracy trade-off
of selective verification, against fixed-strategy, oracle, and rule-based
baselines.
\end{itemize}

To the best of our knowledge, this is the first framework to treat
clarification as a first-class routing action alongside dense, graph, and
hybrid retrieval in a Hybrid GraphRAG pipeline for Vietnamese legal question
answering~\cite{sovrano2020legalkg,zhong2020nlp_law}.

\subsection{Organisation}\label{subsec:organisation}

The remainder of this paper is organised as follows.
Section~\ref{sec:related_work} reviews related work.
Section~\ref{sec:methodology} defines the routing problem, the dataset, and
the overall pipeline. Section~\ref{sec:routing_framework} details the routing
framework. Sections~\ref{sec:eval} and~\ref{sec:results} present the
experimental setup and results. Section~\ref{sec:limitations} discusses
limitations, and Section~\ref{sec:conclusion} concludes.

%% ===== related_work ====================================================
\section{Related Work}\label{sec:related_work}
% ===============================================================

\noindent\textbf{Retrieval-Augmented Generation.}
RAG combines parametric knowledge from neural language models with
non-parametric knowledge retrieved from an external
corpus~\cite{lewis2020rag,guu2020realm}. Dense retrieval is effective for
open-domain factoid QA but degrades on multi-hop evidence
chains~\cite{karpukhin2020dpr,trivedi2023interleaving}. A line of advanced
variants improves retrieval control---Self-RAG through
self-reflection~\cite{asai2023selfrag}, FLARE through active
retrieval~\cite{jiang2023flare}, RAPTOR through hierarchical
summaries~\cite{sarthi2024raptor}, and CRAG through corrective
re-retrieval~\cite{yan2024corrective}. These methods refine \emph{how} dense
retrieval is performed, but they operate within a single evidence mechanism
and do not decide whether a query requires relational traversal or
clarification.

\noindent\textbf{Graph Retrieval and GraphRAG.}
Graph-based retrieval lets evidence follow explicit
relationships~\cite{sun2019pullnet,wang2024kgp}. GraphRAG builds entity
graphs for query-focused summarisation~\cite{edge2024graphrag}; HippoRAG
introduces a neurobiologically inspired graph memory for associative
retrieval~\cite{gutierrez2024hipporag}; LightRAG uses dual-level retrieval
over entity and community knowledge~\cite{he2024lightrag}; and
agent-style methods such as Think-on-Graph~\cite{sun2024thinkongraph} and
GraphReader~\cite{li2024graphreader} let a language model iteratively explore
the graph during reasoning. For legal QA, graph retrieval is especially
relevant because legal documents encode structured relations such as
amendment, replacement, and jurisdiction~\cite{sovrano2020legalkg}. However,
these systems assume graph traversal is always applied; none decides
\emph{when} traversal is warranted versus when a dense lookup would suffice at
lower cost.

\noindent\textbf{Adaptive Routing and Cost-aware Inference.}
Adaptive-RAG selects a retrieval strategy according to question
complexity~\cite{jeong2024adaptive}, while FrugalGPT and RouteLLM route
queries across models of different
cost~\cite{chen2023frugalgpt,ong2024routellm}. These works establish the
cost-aware cascade principle that motivates our design. Their routing target,
however, is the \emph{model} or \emph{retrieval depth}; they do not route
among qualitatively different evidence mechanisms, and they do not treat
clarification as a routable action. Adaptive-RAG is the closest method in
spirit, but it was developed for English open-domain QA and routes among
zero-, single-, and multi-step retrieval over a single dense index, so it is
not directly transferable to a Vietnamese legal corpus that must additionally
choose between dense and graph evidence; we discuss this comparability gap in
Section~\ref{sec:limitations}.

\noindent\textbf{Ambiguity Handling.}
CANARD studies question rewriting for context-dependent
questions~\cite{elgohary2019canard}, and AmbigQA shows that open-domain
questions often admit multiple valid interpretations~\cite{min2020ambigqa}.
These studies treat ambiguity resolution as a stand-alone task applied before
or after retrieval, rather than as one possible outcome of a unified routing
decision.

\noindent\textbf{Vietnamese NLP and Legal QA.}
VnCoreNLP provides Vietnamese word segmentation and
NER~\cite{vu2018vncorenlp}, and PhoBERT provides pretrained Vietnamese
representations~\cite{nguyen2020phobert}. Legal NLP more broadly requires
reasoning over regulatory
relations~\cite{chalkidis2020legalbert,zhong2020nlp_law,guha2023legalbench,louis2024interpretable}.
This body of work supplies the linguistic and legal-domain foundations our
system builds on, but does not address adaptive retrieval routing.

Table~\ref{tab:related_comparison} summarises where the proposed framework
stands relative to these lines of work along the capabilities most relevant to
heterogeneous legal question answering: whether a method selects among
distinct evidence mechanisms rather than a single fixed one, whether that
selection is made \emph{before} retrieval, whether relational (graph) evidence
is supported, whether clarification is a routing option, and whether
conversational context is resolved. No prior line offers all five.

\begin{table}[!htbp]
\centering
\caption{Capabilities of representative retrieval paradigms versus the
proposed framework. \textsf{Mech.}: selects among distinct evidence
mechanisms; \textsf{Pre-ret.}: decides before retrieval; \textsf{Graph}:
supports relational traversal; \textsf{Clarify}: clarification as a routing
option; \textsf{Conv.}: resolves conversational context}\label{tab:related_comparison}
\small
\begin{tabularx}{\columnwidth}{Y ccccc}
\toprule
\textbf{Paradigm} & \textbf{Mech.} & \textbf{Pre-ret.} & \textbf{Graph} & \textbf{Clarify} & \textbf{Conv.} \\
\midrule
Dense RAG~\cite{lewis2020rag,karpukhin2020dpr}          &            &            &            &            &            \\
Advanced RAG~\cite{asai2023selfrag,jiang2023flare}      &            &            &            &            &            \\
GraphRAG~\cite{edge2024graphrag,he2024lightrag}         &            &            & \checkmark &            &            \\
Hybrid RAG~\cite{gao2023modularrag,wang2024kgp}         & \checkmark &            & \checkmark &            &            \\
Adaptive routing~\cite{jeong2024adaptive,ong2024routellm} & \checkmark & \checkmark &          &            &            \\
Ambiguity handling~\cite{min2020ambigqa,elgohary2019canard} &        &            &            & \checkmark & \checkmark \\
\midrule
\textbf{This work} & \checkmark & \checkmark & \checkmark & \checkmark & \checkmark \\
\bottomrule
\end{tabularx}
\end{table}

In summary, despite substantial progress along each of these lines, no prior
system decides among dense retrieval, graph traversal, hybrid reasoning, and
clarification \emph{before} retrieval in a single framework. Advanced RAG
methods improve a single mechanism; GraphRAG methods assume traversal is
always needed; cost-aware routers select models rather than evidence
mechanisms; and ambiguity-handling methods treat clarification in isolation.
Our work differs in placing an explicit, reasoning-aware routing decision at
the centre of the pipeline, and in treating clarification as a first-class
route backed by deterministic conversational context resolution.

%% ===== methodology =====================================================
\section{Methodology}\label{sec:methodology}
% ===============================================================

\subsection{Problem Definition}\label{subsec:problem_definition}

Let $\mathcal{C}=\{D_1,\ldots,D_N\}$ denote the Vietnamese legal corpus of
$N$ documents. From $\mathcal{C}$, the system builds two complementary
knowledge sources: a dense vector index
$\mathcal{I}_v=\{(c_j,\mathbf{e}_j)\}_{j=1}^{M}$ of $M$ legal text chunks,
where $c_j$ is the $j$-th chunk and $\mathbf{e}_j\in\mathbb{R}^d$ is its
$d$-dimensional embedding; and a legal knowledge graph
$\mathcal{G}=(\mathcal{V},\mathcal{E})$ where a node $v\in\mathcal{V}$
represents a legal document, article, semantic concept, agency, or other
entity, and an edge $(u,r,v)\in\mathcal{E}$ represents a typed relation $r$
from source node $u$ to target node $v$, such as \textsc{HAS\_ARTICLE},
\textsc{MENTIONS}, \textsc{AMENDS}, or \textsc{REPLACES}. Table~\ref{tab:notation}
summarises the notation used throughout.

\begin{table}[!htbp]
\centering
\caption{Summary of notation}\label{tab:notation}
\begin{tabularx}{\columnwidth}{L{3cm}Y}
\toprule
\textbf{Symbol} & \textbf{Meaning} \\
\midrule
$\mathcal{C},\,\mathcal{I}_v,\,\mathcal{G}$ & legal corpus, dense index, knowledge graph \\
$q_t,\,H$ & query at turn $t$ and conversation history $H=\{(q_i,a_i)\}_{i=1}^{t-1}$ \\
$\mathcal{D}=\{d_v,d_g,d_h,d_c\}$ & route set: dense, graph, hybrid, clarify \\
$\rho$ & routing function $\rho(q_t,H)\!\to\! d^*\!\in\!\mathcal{D}$ \\
$E_{d^*},E_v,E_g$ & evidence set for route $d^*$; dense and graph evidence \\
$k,\,h$ & dense retrieval depth; maximum traversal depth \\
$\mathbf{f}(q_t,H,r_H)$ & feature vector, $\mathbf{f}\in\mathbb{R}^{n_f}$ \\
$\phi_{\mathrm{amb}}(q_t,H,r_H)$ & ambiguity score, $\in[0,1]$ \\
$\xi(q_t)$ & multi-hop score, $\in\{0,1,2\}$ \\
$r_H$ & history-resolution status (Eq.~\ref{eq:history_states}) \\
$f_1$ & Stage-1 classifier, outputs $(d_1,s_1)$ \\
$s_1$ & Stage-1 confidence (max softmax probability) \\
$\gamma_{\mathrm{conf}},\gamma_{\mathrm{amb}},\gamma_{\mathrm{hop}}$ & binary Stage-2 trigger indicators \\
$\Gamma$ & composite Stage-2 trigger \\
$\tau_{\mathrm{conf}},\tau_{\mathrm{force}},\tau_{\mathrm{clar}}$ & confidence, forcing, and clarify-override thresholds \\
$f_2$ & Stage-2 LLM verifier \\
$a_t,\,c_t$ & generated answer; clarification question \\
\bottomrule
\end{tabularx}
\end{table}

At inference time, the input is the current query $q_t$ at turn $t\geq 1$,
together with optional conversation history $H=\{(q_i,a_i)\}_{i=1}^{t-1}$
(empty when $t=1$). The router selects one action from the route set
\begin{equation}
\mathcal{D}=\{d_v,\,d_g,\,d_h,\,d_c\},
\label{eq:route_set}
\end{equation}
where $d_v$, $d_g$, $d_h$, $d_c$ denote \texttt{dense\_retrieval},
\texttt{graph\_traversal}, \texttt{hybrid\_reasoning}, and \texttt{clarify},
respectively. The router is modelled as a function $\rho$ mapping the input
pair to a route:
\begin{equation}
d^* = \rho(q_t,H),\qquad d^*\in\mathcal{D}.
\label{eq:routing_function}
\end{equation}
For the three retrieval routes, the evidence set $E_{d^*}$ is determined by
the selected route $d^*$:
\begin{equation}
E_{d^*}=
\begin{cases}
\mathrm{TopK}(q_t,\mathcal{I}_v,k) & \text{if } d^*=d_v,\\[3pt]
\mathrm{Traverse}(q_t,\mathcal{G},h) & \text{if } d^*=d_g,\\[3pt]
E_v\cup E_g & \text{if } d^*=d_h,
\end{cases}
\label{eq:evidence_set}
\end{equation}
where $E_v\triangleq\mathrm{TopK}(q_t,\mathcal{I}_v,k)$ is the set of the $k$
chunks most similar to $q_t$ in $\mathcal{I}_v$;
$E_g\triangleq\mathrm{Traverse}(q_t,\mathcal{G},h)$ is the evidence obtained
by linking entities in $q_t$ to nodes of $\mathcal{G}$ and traversing up to a
maximum depth $h$; $k\in\mathbb{N}^{*}$ is the dense retrieval depth (number
of top-$k$ chunks returned) and $h\in\mathbb{N}^{*}$ is the maximum traversal
depth. The clarify route $d_c$ produces no evidence set and therefore does
not appear in~\eqref{eq:evidence_set}. The generator then produces
$a_t=\mathrm{Generate}(q_t,H,E_{d^*})$. If $d^*=d_c$, no evidence is
retrieved; instead the system returns a clarification question
$c_t=\mathrm{Clarify}(q_t,H)$.

\subsection{Dataset}\label{subsec:dataset}

The evaluation uses a self-constructed Vietnamese legal QA routing dataset
derived from a single source: the Vietnamese electronic legal codification
system (the ``Vietnamese Legal Code Portal''), flattened into article-level
records.
Question--answer pairs were generated from sampled legal articles using the
Gemini API. For each text snippet, the model was prompted with templates
specific to the three evidence structures to synthesize realistic user
queries, and the exact corresponding span in the article was extracted as the
ground-truth answer. A subset of the generated pairs was then manually
reviewed to ensure legal soundness and prevent hallucinated facts.
QA pairs are linked to supporting articles through an \texttt{article\_key}
(format \texttt{DOC::ARTICLE\_ID}) and assigned routing labels based on
evidence structure: \texttt{dense\_retrieval} for single-article questions,
\texttt{graph\_traversal} for multi-article reasoning within one document,
and \texttt{hybrid\_reasoning} for evidence spanning multiple documents.
Low-quality samples are removed by deduplication, length filtering, and
placeholder-answer filtering. To reduce template-induced lexical leakage in
the routing features, the surface forms of the generated queries are
additionally paraphrased while preserving their underlying evidence
structure, so that the routing classifier must rely on genuine structural
signals rather than memorised template wording.

Table~\ref{tab:dataset} summarises the dataset splits. The dataset remains
skewed toward dense retrieval, which accounts for half of the samples
($50\%$; Fig.~\ref{fig:dataset_bar} shows the test-split distribution),
motivating Macro-F1 reporting alongside accuracy. The strict split
is used for retrieval-route evaluation; the \texttt{clarify} route is
evaluated separately on a $234$-query ambiguity benchmark ($156$ ambiguous,
$78$ answerable controls) and a $160$-query conversation stress test.
The $234$-query ambiguity dataset was constructed by systematically dropping
essential context (such as subjects or conditions) from well-posed queries to
form $156$ ambiguous samples, which were paired with $78$ unambiguous
controls authored by legal experts. The $160$-query conversation set was
synthesized to simulate realistic multi-turn dialogues encompassing topic
shifts and contextual references, followed by a human validation pass to
guarantee contextual coherence.

\begin{table}[!htbp]
\centering
\caption{Vietnamese legal QA routing dataset split}\label{tab:dataset}
\begin{tabular}{lrrrr}
\toprule
\textbf{Split} & \textbf{Dense} & \textbf{Graph} & \textbf{Hybrid} & \textbf{Total} \\
\midrule
Train & 400 & 200 & 200 & 800 \\
Dev   & 50  & 25  & 25  & 100 \\
Test  & 300 & 150 & 150 & 600 \\
\midrule
\textbf{Total} & 750 & 375 & 375 & 1{,}500 \\
\bottomrule
\end{tabular}
\end{table}

\begin{figure}[!htbp]
\centering
\includegraphics[width=0.77\columnwidth]{dataset_bar.pdf}
\caption{Route-label distribution of the strict $600$-query evaluation set:
$300$ dense-retrieval, $150$ graph-traversal, and $150$ hybrid-reasoning
queries}\label{fig:dataset_bar}
\end{figure}

\subsection{System Overview}\label{subsec:system_overview}

\begin{figure}[!htbp]
\centering
\includegraphics[width=\columnwidth]{system_architecture.pdf}
\caption{High-level data flow of the proposed Hybrid GraphRAG system. The
Vietnamese legal corpus $\mathcal{C}$ is processed offline into a dense
vector index $\mathcal{I}_v$ (ChromaDB) and a legal knowledge graph
$\mathcal{G}$ (Neo4j). At query time the two-stage adaptive router selects
evidence from $\mathcal{I}_v$, $\mathcal{G}$, or both, or returns a
clarification question $c_t$. Box colours denote input, processing, storage,
routing, and output stages}\label{fig:system_architecture}
\end{figure}

\begin{figure*}[!htbp]
\centering
\includegraphics[width=\textwidth]{system_pipeline.pdf}
\caption{Reasoning-aware adaptive routing pipeline and illustrative routing
outcomes. \textbf{a} The pipeline processes $(q_t,H)$ through feature
extraction, Stage~1 routing, trigger checking, optional Stage~2 LLM
verification, and route-controlled retrieval or clarification. \textbf{b}
Concrete examples for the four routing actions: dense ($d_v$) for direct
lookup, graph ($d_g$) for amendment relations, hybrid ($d_h$) for
cross-document synthesis, and clarify ($d_c$) for an underspecified
referent}\label{fig:system_pipeline}
\end{figure*}

Figure~\ref{fig:system_architecture} shows the overall data flow. In the
offline phase, $\mathcal{C}$ is indexed once into $\mathcal{I}_v$ and
$\mathcal{G}$. In the online phase, every $(q_t,H)$ is routed against these
fixed sources. Figure~\ref{fig:system_pipeline} illustrates both the routing
pipeline and representative examples of the four routing actions. The final
route $d^*$ controls evidence retrieval, graph reasoning, hybrid evidence
aggregation, or clarification generation.

\subsection{Implementation Details}\label{subsec:implementation}

The Chroma vector store is persisted with chunks embedded by
\texttt{microsoft/Harrier-OSS-v1-0.6B} ($256$-token chunks, $32$-token
overlap, top-$k=8$). The graph backend uses Neo4j with full-text indexes over
legal documents, articles, and vector chunks; graph retrieval uses
Text-to-Cypher querying via an OpenAI-compatible endpoint that generates
Cypher queries against $\mathcal{G}$. Answer generation and the Stage-2
verifier are served by the same local vLLM endpoint running a
\texttt{Qwen3-32B-AWQ} model; decoding uses greedy search (temperature~$0$)
with a fixed seed, so that answers are deterministic and any run-to-run
variation in the reported metrics comes only from resampling queries
(Section~\ref{subsec:metrics}).
Answers are generated with a citation-constrained prompt that
requires verbatim quotation of only the clause or point directly answering
the question (quoting an entire article when unnecessary is treated as an
error), permits a short direct answer before the supporting quotation, and
prescribes a single fixed sentence when the retrieved evidence contains no
relevant provision; the full system prompt is reproduced in
Appendix~\ref{app:repro}. Stage~2 is prompted for a structured JSON route
prediction (route decision, confidence, an optional clarification question,
and a reasoning trace).
Vietnamese named-entity recognition uses
\texttt{NlpHUST/ner-vietnamese-electra-base}~\cite{clark2020electra}. The
Stage~1 router checkpoint is stored separately from the service layer, so
that the router can be retrained without redeploying the whole system. Full
hyperparameters, the complete feature list, the ambiguity-score definition,
the history-resolver rules, and the Stage-2 prompt are given in
Appendix~\ref{app:repro}.

%% ===== framework =======================================================
\section{Reasoning-aware Adaptive Routing Framework}\label{sec:routing_framework}
% ===============================================================

\subsection{Legal Knowledge Base}\label{subsec:legal_kb}

The vector index $\mathcal{I}_v$ contains $68{,}835$ embedded chunks
appropriate for direct questions whose answers are stated in legal text. The
graph $\mathcal{G}$ stores explicit typed relations among legal documents,
articles, and concepts. Tables~\ref{tab:graph_stats}
and~\ref{tab:graph_relations} report the statistics of the graph built from
the national legal codification, in which article nodes are promoted to
\textsc{LegalArticle} to enable article-level evidence retrieval.

\begin{table}[!htbp]
\centering
\caption{Vietnamese legal graph $\mathcal{G}$ statistics}\label{tab:graph_stats}
\begin{tabular}{lr}
\toprule
\textbf{Graph item} & \textbf{Count} \\
\midrule
Nodes $|\mathcal{V}|$ & 804{,}322 \\
Edges $|\mathcal{E}|$ & 1{,}968{,}621 \\
Legal documents & 4{,}598 \\
Legal articles (\textsc{LegalArticle}) & 68{,}570 \\
Semantic concepts & 392{,}592 \\
Other entities & 338{,}562 \\
\bottomrule
\end{tabular}
\end{table}

\begin{table}[!htbp]
\centering
\caption{Main relation types in the legal graph, ordered by count}\label{tab:graph_relations}
\begin{tabular}{lr}
\toprule
\textbf{Relation type} & \textbf{Count} \\
\midrule
\textsc{MENTIONS} & 1{,}033{,}341 \\
\textsc{HAS\_CONDITION} & 202{,}298 \\
\textsc{OBLIGATES} & 128{,}768 \\
\textsc{PART\_OF\_PROCEDURE} & 102{,}004 \\
\textsc{REGULATES} & 75{,}137 \\
\bottomrule
\end{tabular}
\end{table}

The graph is large and densely linked at the semantic level through relations
such as \textsc{MENTIONS} ($1{,}033{,}341$ edges) and \textsc{HAS\_CONDITION}
($202{,}298$), whereas explicit cross-document legal-effect relations such as
\textsc{AMENDS} ($5{,}944$ edges) and \textsc{REPLACES} ($4{,}189$) remain
comparatively sparse. This asymmetry directly motivates adaptive routing:
graph traversal is valuable for relation-heavy queries but unnecessarily
costly for simple article-level lookups.

\subsection{Reasoning-aware Feature Extraction}\label{subsec:features}

For each input $(q_t,H)$, the router first runs the history resolver
(Eq.~\ref{eq:history_resolver}) and then extracts a feature vector
$\mathbf{f}(q_t,H,r_H)\in\mathbb{R}^{n_f}$ that blends lexical signals with
reasoning-aware metrics, together with two auxiliary scores: the ambiguity
score $\phi_{\mathrm{amb}}(q_t,H,r_H)\in[0,1]$ and the multi-hop score
$\xi(q_t)\in\{0,1,2\}$.
Here $n_f=27$, comprising $16$ lexical surface features produced by a
Vietnamese legal feature extractor and $11$ reasoning-aware complexity
features produced by a query-complexity analyser; the complete list, in the
canonical vector order used for both training and inference, is given in
Appendix~\ref{app:repro}, Table~\ref{tab:features}. Using an explicit feature
vector rather than a hidden query embedding keeps Stage~1 both interpretable
and cheap to evaluate.
The Stage~1 classifier relies most heavily on a small set of structural
features: the \textbf{legal reference count}, which tracks how many distinct
legal documents are mentioned; the \textbf{query length in words}, as
cross-document synthesis questions tend to be longer;
\textbf{cross-document signals}, reflecting explicit synthesis indicators;
and the \textbf{article reference count}. Additional reasoning
metrics---multi-hop connectors, sub-question count, procedural depth, and a
discrete complexity level---provide finer-grained separation for harder
cases, while question-type cues (``what'', ``when'', ``how much'') and
legal-relation cues (``amend'', ``replace'', ``takes effect from'') further
inform the decision.

\paragraph{Multi-hop score.}
Let $\mathcal{K}_{\mathrm{rel}}$ be the fixed lexicon of legal-relation cue
phrases (e.g.\ ``amend'', ``replace'', ``repeal'', ``provide guidance'',
``pursuant to''), listed in full in Appendix~\ref{app:repro}. The multi-hop
score counts how many distinct cues occur in the query, capped at two:
\begin{equation}
\xi(q_t)=\min\!\Bigl(2,\;\bigl|\{\,\kappa\in\mathcal{K}_{\mathrm{rel}} : \kappa\ \text{occurs in}\ q_t\,\}\bigr|\Bigr).
\label{eq:multihop_score}
\end{equation}
Thus $\xi(q_t)=0$ means no relation signal (likely a direct lookup),
$\xi(q_t)=1$ means a single relation cue, and $\xi(q_t)=2$ means two or more
co-occurring cues, indicating the strongest multi-hop demand.

\paragraph{Ambiguity score.}
The ambiguity score is produced by a deterministic detector that evaluates a
set of $m=11$ binary ambiguity indicators
$\{\psi_1(q_t,H,r_H),\ldots,\psi_m(q_t,H,r_H)\}$---for example, an
unresolved demonstrative or pronoun (e.g.\ ``that'', ``this'', ``it'')
whose referent is not fixed by $r_H$, a generic legal need with no concrete
legal-object anchor, or an elliptical follow-up with no usable history---each
carrying a calibrated weight $w_i\in[0,1]$. The score is the strongest
triggered signal, gated by the history-resolution status:
\begin{equation}
\phi_{\mathrm{amb}}(q_t,H,r_H)=
\begin{cases}
0 & \text{if } r_H=\texttt{resolved},\\[3pt]
\displaystyle\max_{1\le i\le m}\; w_i\,\psi_i(q_t,H,r_H) & \text{otherwise},
\end{cases}
\label{eq:ambiguity_score}
\end{equation}
so that a referent recovered from history suppresses the ambiguity signal
regardless of surface cues. Taking the maximum rather than a sum reflects the
weight calibration: a single strong indicator (e.g.\ an unresolved pronoun
with empty history, $w=0.90$) is sufficient to exceed the trigger threshold
$\tau_{\mathrm{force}}=0.60$, whereas weak indicators ($w\le 0.30$) cannot
force verification on their own. The full indicator set and weights
$\{(\psi_i,w_i)\}$, including two graded indicator families whose scores
scale with match counts, are specified in Appendix~\ref{app:repro},
Table~\ref{tab:ambiguity_features}.

\paragraph{History resolver.}
A deterministic history resolver runs before ambiguity scoring,
\begin{equation}
r_H = \mathrm{Resolve}(q_t,H),
\label{eq:history_resolver}
\end{equation}
returning a referent-resolution status
\begin{equation}
r_H\in\{\texttt{resolved},\,\texttt{no\_history},\,\texttt{irrelevant\_history},\,\texttt{conflicting\_history},\,\texttt{not\_needed}\}.
\label{eq:history_states}
\end{equation}
If a phrase such as ``that document'' has a unique document referent in
$H$, the resolved referent can unblock graph traversal; otherwise the same
surface form remains ambiguous. The status $r_H$ is used both as a Stage~1
input feature and as one of the Stage~2 trigger conditions; the resolution
rules are listed in Appendix~\ref{app:repro}.

\subsection{Stage~1: XGBoost Router}\label{subsec:stage1}

Stage~1 is an XGBoost multi-class classifier~\cite{chen2016xgboost} trained on
the $800$ routing-labelled training queries augmented with $60$ synthetic
\texttt{clarify} samples drawn from ambiguous-query templates, since the QA
dataset itself contains no gold \texttt{clarify} label:
\begin{equation}
(d_1,s_1) = f_1\bigl(\mathbf{f}(q_t,H,r_H)\bigr),
\label{eq:stage1}
\end{equation}
where $d_1\in\mathcal{D}=\{d_v,d_g,d_h,d_c\}$ is the initial route and
$s_1\in[0,1]$ is the confidence, taken as the maximum class probability of
the softmax output. Because the \texttt{clarify} class is only weakly
supervised by synthetic templates, a Stage~1 prediction $d_1=d_c$ is treated
as a proposal rather than a decision: whenever $d_1=d_c$, Stage~2 is invoked
to adjudicate, so clarification is never emitted on the strength of the
weakly supervised class alone.
XGBoost is chosen for three reasons: its
features are mostly discrete counts and binary indicators that
gradient-boosted trees handle without elaborate normalisation; inference over
$n_f$ features is orders of magnitude cheaper than one LLM pass, meeting the
latency-aware requirement of a stage that processes all traffic; and tree
models expose feature importances that support interpretability, which is
valuable in the legal domain. Empirically, the classifier forward pass takes
$\approx\!8.9$~ms per query on CPU (the full Stage~1 path, including feature
extraction and history resolution, is $\approx\!33.7$~ms; see
Section~\ref{subsec:rq5}), versus $\approx\!53.3$~ms for a
PhoBERT~\cite{nguyen2020phobert} classifier on GPU, making Stage~1 suitable as
a real-time gateway and motivating a future three-level cascade
XGBoost~$\rightarrow$~PhoBERT~$\rightarrow$~LLM following
FrugalGPT~\cite{chen2023frugalgpt}.

\subsection{Stage~2: Selective LLM Verifier}\label{subsec:stage2}

Stage~2 is an LLM verifier invoked selectively. Three binary trigger
indicators, each taking values in $\{0,1\}$, are defined as
\begin{align}
\gamma_{\mathrm{conf}} &= \mathbb{I}[\,s_1 < \tau_{\mathrm{conf}}\,],
    \label{eq:gamma_conf}\\
\gamma_{\mathrm{amb}}  &= \mathbb{I}[\,\phi_{\mathrm{amb}}(q_t,H,r_H) \geq \tau_{\mathrm{force}}\,],
    \label{eq:gamma_amb}\\
\gamma_{\mathrm{hop}}  &= \mathbb{I}[\,\xi(q_t) \geq 2\,],
    \label{eq:gamma_hop}
\end{align}
where $\mathbb{I}[\cdot]$ is the indicator function (equal to $1$ when its
argument is true and $0$ otherwise). Thus $\gamma_{\mathrm{conf}}$ fires when
Stage~1 confidence is below $\tau_{\mathrm{conf}}$, $\gamma_{\mathrm{amb}}$
fires when the ambiguity score reaches the forcing threshold
$\tau_{\mathrm{force}}$, and $\gamma_{\mathrm{hop}}$ fires at the strongest
multi-hop level ($\xi=2$). With the convention $a\vee b=\max(a,b)$ on
$\{0,1\}$, the composite trigger is the logical disjunction of the three
indicators together with the conflicting-history condition,
\begin{equation}
\Gamma(q_t,H,r_H) = \gamma_{\mathrm{conf}}\vee\gamma_{\mathrm{amb}}\vee\gamma_{\mathrm{hop}}\vee\mathbb{I}[\,r_H=\texttt{conflicting\_history}\,]\;\in\;\{0,1\},
\label{eq:trigger_function}
\end{equation}
so that $r_H$ enters the trigger explicitly.

When $\Gamma=0$, Stage~1 is trusted and $d^*=d_1$. When $\Gamma=1$, the
verifier is queried and returns a proposed route $\tilde{d}\in\mathcal{D}$
together with, when $\tilde{d}=d_c$, a clarify confidence
$p_c\in[0,1]$ (the model's estimated probability that the query cannot be
safely grounded as posed). A clarification is emitted only if the verifier is
sufficiently confident; otherwise the proposed retrieval route is kept. The
clarify-override threshold $\tau_{\mathrm{clar}}$ therefore acts on the
verifier output, and the final route is
\begin{equation}
d^* =
\begin{cases}
d_1 & \text{if } \Gamma(q_t,H,r_H)=0,\\[3pt]
d_c & \text{if } \Gamma=1,\ \tilde{d}=d_c,\ \text{and } p_c\ge\tau_{\mathrm{clar}},\\[3pt]
\tilde{d} & \text{otherwise } (\Gamma=1,\ \tilde{d}\in\{d_v,d_g,d_h\}\ \text{or } p_c<\tau_{\mathrm{clar}}),
\end{cases}
\label{eq:final_route}
\end{equation}
where in the last case a below-threshold clarify proposal is demoted to the
verifier's best retrieval route. Default thresholds are
$\tau_{\mathrm{conf}}=0.50$, $\tau_{\mathrm{force}}=0.60$, and
$\tau_{\mathrm{clar}}=0.80$; their effect is analysed in
Section~\ref{subsec:rq4}. In words, Stage~2 is invoked on low-confidence
Stage~1 decisions, on high ambiguity risk, on strong multi-hop signals, or on
conflicting history; a \texttt{resolved} history unblocks retrieval and
suppresses false clarification, whereas missing, irrelevant, or conflicting
history forces verification or clarification.

Algorithm~\ref{alg:router} summarises the procedure. The deployed
implementation additionally uses two heuristics: a dense fast path that skips
Stage~2 for high-confidence, low-ambiguity dense queries, and a
reasoning-signal check that escalates queries with strong legal features even
at moderate Stage~1 confidence.

\begin{algorithm}[!htbp]
\caption{Two-stage adaptive router}\label{alg:router}
\begin{algorithmic}[1]
\Require query $q_t$, history $H$, thresholds $\tau_{\mathrm{conf}},\tau_{\mathrm{force}},\tau_{\mathrm{clar}}$
\Ensure final route $d^*\in\mathcal{D}$
\State $r_H \gets \mathrm{Resolve}(q_t,H)$
\State $(\mathbf{f},\phi_{\mathrm{amb}},\xi) \gets \mathrm{ExtractFeatures}(q_t,H,r_H)$
\State $(d_1,s_1) \gets f_1(\mathbf{f})$
\State $\gamma_{\mathrm{conf}}\gets\mathbb{I}[s_1<\tau_{\mathrm{conf}}]$;\ \ $\gamma_{\mathrm{amb}}\gets\mathbb{I}[\phi_{\mathrm{amb}}\ge\tau_{\mathrm{force}}]$;\ \ $\gamma_{\mathrm{hop}}\gets\mathbb{I}[\xi\ge 2]$
\State $\Gamma \gets \gamma_{\mathrm{conf}}\vee\gamma_{\mathrm{amb}}\vee\gamma_{\mathrm{hop}}\vee\mathbb{I}[r_H=\texttt{conflicting\_history}]$
\If{$\Gamma=0$}
  \State \Return $d_1$
\EndIf
\State $(\tilde{d},p_c) \gets f_2(q_t,H,d_1,s_1,\mathbf{f})$
\If{$\tilde{d}=d_c$ \textbf{and} $p_c\ge\tau_{\mathrm{clar}}$}
  \State \Return $d_c$
\Else
  \State \Return $\tilde{d}$ \Comment{demote below-threshold clarify to best retrieval route}
\EndIf
\end{algorithmic}
\end{algorithm}

\noindent\textbf{Running example.}
For $q_1=$ ``What are the conditions for issuing a building permit?'',
$\xi(q_1)=0$ and Stage~1 predicts $d_v$ with high confidence; Stage~2 is
bypassed. For $q_2=$ ``Which decree amends the administrative-penalty rules
on land?'', the amendment cue (``amend'') raises relational signals, so the
router selects $d_g$ or $d_h$. For $q_3=$ ``Is that document still in
effect?'' with $H=\emptyset$, the referent is unresolved,
$\phi_{\mathrm{amb}}$ is high, $\Gamma=1$, and the verifier returns
$\tilde{d}=d_c$ with $p_c\ge\tau_{\mathrm{clar}}$.

\subsection{Route Execution and Fallback Policy}\label{subsec:route_execution}

For $d_v$, the system retrieves top-$k$ chunks from $\mathcal{I}_v$. For
$d_g$, it links query entities to graph nodes and traverses $\mathcal{G}$. For
$d_h$, it constructs $E_{d^*}=E_v\cup E_g$. For $d_c$, it returns a
clarification question without retrieval. A defensive fallback policy applies
when primary retrieval is uninformative: a failed graph traversal or hybrid
synthesis falls back to dense retrieval; an uninformative dense lookup expands
$k$ once and retries; and a clarify route, if unanswered, keeps the
conversation in a waiting state rather than guessing the referent.

%% ===== experimental_setup ==============================================
\section{Experimental Setup}\label{sec:eval}
% ===============================================================

\subsection{Research Questions}\label{subsec:rqs}

We organise the evaluation around five research questions that separate the
two claims of the paper---that routing among evidence mechanisms is both
learnable and useful, and that a selective second stage adds value beyond
route correction.

\begin{description}
\item[RQ1 (End-to-end quality).] Does adaptive routing improve answer quality
over fixed-strategy and oracle baselines, and does Stage~2 contribute beyond
mere routing correction?
\item[RQ2 (Routing accuracy).] How accurately does the Stage~1 classifier
select routes under cross-validation, compared with rule-based and learned
baselines?
\item[RQ3 (Ambiguity handling).] Can the system avoid unsafe retrieval by
clarifying when the legal target is underspecified?
\item[RQ4 (Threshold ablation).] How do the key thresholds affect routing
accuracy and the Stage~2 trigger rate?
\item[RQ5 (Cost--quality).] Does selective Stage~2 invocation provide a better
trade-off than always-on or always-off verification?
\end{description}

\subsection{Systems and Baselines}\label{subsec:systems}

Table~\ref{tab:systems} defines all compared systems. All end-to-end
experiments run on the \texttt{phapdien\_strict} test set of $600$ queries,
using the single national codification corpus with one shared retrieval and
generation configuration, so that any observed difference reflects routing
decisions rather than infrastructure. The Oracle Router dispatches every query
using its gold route label under identical retrieval and generation settings;
comparing the oracle against the Two-stage Hybrid therefore quantifies how
much answer quality is attributable to Stage~2 reasoning beyond correct
routing.

\begin{table}[!htbp]
\centering
\caption{Systems compared in this work}\label{tab:systems}
\begin{tabularx}{\columnwidth}{L{2.7cm}Y}
\toprule
\textbf{System} & \textbf{Description} \\
\midrule
Pure Vector & Always dense retrieval from Chroma \\
Pure Graph & Always graph traversal in Neo4j (Text-to-Cypher) \\
Pure Hybrid & Always combines dense and graph evidence \\
Single-stage Router & XGBoost Stage~1 only; no LLM verifier \\
Two-stage Hybrid & Stage~1 + selective Stage~2 (full system) \\
Always-on Stage~2 & Stage~1 + unconditional Stage~2 verification (trigger-policy ablation) \\
Oracle Router & Gold route labels; retrieval only, no Stage~2 \\
Oracle + Stage~2 & Gold route labels \emph{with} Stage~2 reasoning (CoT-isolation ablation) \\
\bottomrule
\end{tabularx}
\end{table}

\subsection{Evaluation Metrics}\label{subsec:metrics}

\paragraph{Answer quality.}
For answer quality we report Exact Match (EM) and token-level F1 computed by a
single \texttt{evaluate\_prediction(gold\_context, answer)} routine shared by
all systems; by convention F1$\,=0$ when precision and recall are both zero.
Because token-F1 against a gold context rewards lexical overlap and can
penalise a correct answer phrased differently from the reference---a bias that
disadvantages systems whose evidence is assembled from structured graph
context rather than copied from a passage---we additionally report a semantic
metric to complement lexical F1.
We report BERTScore-F1~\cite{zhang2020bertscore} computed with the
Vietnamese encoder \texttt{vinai/phobert-base}~\cite{nguyen2020phobert} (first
$12$ layers, \texttt{lang="vi"}) as this semantic complement, so that answers
scored as semantically faithful but lexically divergent from the gold span
are credited rather than penalised.
Hit@$k$ measures retrieval coverage (the gold source
appears within the top $k$ results) over the source-labelled queries. Because
the dense index, the graph store, and the gold annotations historically used
four incompatible identifier formats, both retrieved and gold references are
first mapped to a canonical \texttt{law::\allowbreak document::\allowbreak
article::\allowbreak number} form by a deterministic normaliser before
matching (strict mode: document code \emph{and} article number must agree);
references the normaliser cannot resolve are counted as misses and flagged,
rather than silently matched (Appendix~\ref{app:repro}).

\paragraph{Route classification.}
For route classification we report accuracy, per-class F1, and Macro-F1 over
the class set $\mathcal{R}=\{d_v,d_g,d_h\}$; Macro-F1 is emphasised because the
dataset is dominated by the dense class~\cite{powers2011evaluation}. To avoid
reporting an over-optimistic single-split figure, the headline routing result
is obtained by \emph{stratified five-fold cross-validation}, and we report the
mean and standard deviation across folds.

\paragraph{Clarification.}
For clarification we report precision, recall, and F1 of the \texttt{clarify}
decision on the ambiguity benchmarks, computed against the binary
ambiguous/answerable gold label.

\paragraph{Cost.}
For cost we report average and 95th-percentile (p95) latency, the Stage~2
trigger rate, and the graph-to-dense fallback rate. To attribute latency to its
sources rather than to the pipeline as a whole, we additionally decompose
end-to-end latency into a routing component (Stage~1, optional Stage~2) and a
retrieval-plus-generation component (Section~\ref{subsec:rq5}).

\paragraph{Statistical reporting.}
Because all end-to-end scores are means over a finite test set ($N=600$),
point estimates alone can overstate the separation between systems. We
therefore accompany the headline F1 and Hit@$1$ figures with $95\%$
confidence intervals obtained by bootstrap resampling of the per-query scores
($B=1000$ resamples, fixed seed~$42$), and we regard two systems as
distinguishable only when their intervals do not overlap.

%% ===== results =========================================================
\section{Results and Discussion}\label{sec:results}
% ===============================================================

\subsection{RQ1: End-to-end Answer Quality}\label{subsec:rq1}

Since the ultimate purpose of routing is to produce better answers, we first
ask whether adaptive routing improves end-to-end quality over fixed-strategy
baselines, and whether Stage~2 contributes beyond correcting route selection.
Two ablations isolate the mechanisms: an Oracle Router that dispatches every
query using its gold route label without Stage~2, and an
\textit{Oracle~+~Stage~2} configuration that holds routing at gold labels
while enabling the verifier, so that any difference between the two is
attributable to Stage~2 reasoning alone. The trigger policy is examined
through its two endpoints: the selective Two-stage system and an Always-on
variant that invokes the verifier on every query.
Table~\ref{tab:end_to_end} reports results on the strict $600$-query
benchmark with bootstrap $95\%$ confidence intervals ($B=1000$, seed~$42$),
now including the semantic BERT-F1 complement and canonical Hit@$1$ after the
retrieval-stage identifier fix; Fig.~\ref{fig:f1_bar} visualises Answer~F1.

\begin{table*}[!htbp]
\centering
\caption{End-to-end results on the strict $600$-query Vietnamese legal QA
benchmark. Route~Acc is the agreement between the \emph{executed} route and
the gold label, measured after Stage-2 overrides, clarification decisions,
and execution-time fallbacks. F1 and EM are token-level scores from
\texttt{evaluate\_prediction}, with bootstrap $95\%$ confidence intervals in
brackets; BERT-F1 is the semantic complement
(Section~\ref{subsec:metrics}). Hit@$1$ uses canonical identifier matching in
strict mode. S2 is the Stage~2 trigger rate. Bold marks the best value per
column, excluding the two oracle rows}\label{tab:end_to_end}
\small
\resizebox{\textwidth}{!}{%
\begin{tabular}{lrlrcrrr}
\toprule
\textbf{System} & \textbf{Route Acc} & \textbf{F1 {\scriptsize[95\% CI]}} & \textbf{EM} & \textbf{BERT-F1} & \textbf{Hit@1} & \textbf{Lat (ms)} & \textbf{S2} \\
\midrule
Pure Vector  & 0.500 & 0.663 {\scriptsize[0.635, 0.690]} & 0.152 & 0.836 & \textbf{0.727} & 10{,}914 & -- \\
Pure Graph   & 0.250 & 0.633 {\scriptsize[0.607, 0.659]} & \textbf{0.335} & 0.835 & 0.545 & \textbf{6{,}034} & -- \\
Pure Hybrid  & 0.250 & 0.647 {\scriptsize[0.623, 0.671]} & 0.153 & \textbf{0.865} & 0.523 & 7{,}375 & -- \\
Single-stage & \textbf{0.737} & 0.457 {\scriptsize[0.430, 0.486]} & 0.130 & 0.735 & 0.215 & 7{,}262 & 0.0\% \\
Two-stage (ours) & 0.623 & 0.602 {\scriptsize[0.576, 0.628]} & 0.243 & 0.815 & 0.508 & 9{,}799 & 79.2\% \\
Always-on Stage~2 & 0.455 & \textbf{0.682} {\scriptsize[0.656, 0.706]} & \textbf{0.335} & 0.861 & 0.588 & 10{,}270 & 100.0\% \\
\midrule
Oracle       & 1.000 & 0.524 {\scriptsize[0.497, 0.550]} & 0.113 & 0.771 & 0.415 & 7{,}548 & -- \\
Oracle + Stage~2 & 0.637 & 0.615 {\scriptsize[0.589, 0.641]} & 0.257 & 0.823 & 0.510 & 9{,}920 & 79.2\% \\
\bottomrule
\end{tabular}%
}
\end{table*}

\begin{figure}[!htbp]
\centering
\includegraphics[width=\columnwidth]{f1_bar.pdf}
\caption{Answer~F1 on the strict $600$-query benchmark with bootstrap $95\%$
confidence intervals. Always-on Stage~2 attains the highest F1 ($0.682$),
statistically indistinguishable from the strongest fixed-strategy baseline
(Pure Vector, $0.663$); the selective Two-stage Router reaches $0.602$ and
outperforms the Oracle ($0.524$), while the Single-stage Router ($0.457$) is
the weakest, underscoring the role of Stage~2}\label{fig:f1_bar}
\end{figure}

We make the following observations.

\noindent\textbf{Stage~2 reasoning causally improves answers, on both
metrics.} The \textit{Oracle~+~Stage~2} configuration reaches F1~$=0.615$
{\scriptsize[$0.589, 0.641$]} against F1~$=0.524$
{\scriptsize[$0.497, 0.550$]} for the plain Oracle---a $+0.091$ gain with
\emph{non-overlapping} confidence intervals under \emph{identical} gold
routing---and lifts the semantic BERT-F1 by $+0.052$ ($0.823$ vs $0.771$).
Because routing is held fixed, both differences are attributable to
Stage~2's explicit reasoning over query intent and legal constraints during
verification and generation, confirming the chain-of-thought interpretation
that the Two-stage$\,>\,$Oracle contrast ($0.602$ vs $0.524$, $+0.078$, also
CI-separated) could only suggest.

\noindent\textbf{Stage~2 accounts for the routing-to-quality gain.} The gap
between the Single-stage Router (F1~$=0.457$) and the Two-stage Hybrid
(F1~$=0.602$) is $+0.145$ with non-overlapping intervals, and grows to
$+0.225$ for the Always-on variant. Because all three systems share an
identical Stage~1, these gaps are attributable entirely to Stage~2
verification and reasoning.

\noindent\textbf{Better answers, not better route agreement.} Executed-route
accuracy \emph{falls} monotonically as the verifier intervenes more: from
$0.737$ (Single-stage) to $0.623$ (selective) to $0.455$
(Always-on)---and even Oracle~+~Stage~2 retains only $0.637$ agreement with
the gold labels because the verifier overrides them---yet F1 rises at every
step, peaking exactly where agreement is lowest. Stage~2's value therefore
lies in reasoning during answer generation rather than in reproducing the
gold routing taxonomy, which is precisely the chain-of-thought behaviour
claimed.

\noindent\textbf{Always-on verification is the quality-optimal operating
point.} Applied to every query, the framework reaches F1~$=0.682$
{\scriptsize[$0.656, 0.706$]}: statistically indistinguishable from the
strongest fixed baseline (Pure Vector, $0.663$
{\scriptsize[$0.635, 0.690$]}) and nominally above it, tied with Pure Graph
for the best exact match ($0.335$, $2.2\times$ Pure Vector's $0.152$), and
within $0.004$ of the best semantic score ($0.861$ vs Pure Hybrid's
$0.865$). No fixed strategy matches this profile across all three answer
metrics simultaneously, and only the routed systems can additionally refuse
to guess by clarifying (Section~\ref{subsec:rq3}).

\noindent\textbf{Selective triggering trades measurable quality for a small
latency saving.} The selective Two-stage system ($0.602$
{\scriptsize[$0.576, 0.628$]}) sits \emph{significantly below} both the
Always-on variant ($-0.080$, non-overlapping intervals) and Pure Vector
($-0.061$, non-overlapping) on aggregate token-F1, while saving only
$471$~ms ($4.6\%$) of average latency relative to Always-on
(Section~\ref{subsec:rq5}). On this generation-dominated deployment the
verifier should therefore run on every query; the selective policy becomes
attractive only when verifier cost dominates, as quantified in RQ4/RQ5. The
selective system still exceeds Pure Vector's exact match by $60\%$ ($0.243$
vs $0.152$).

\noindent\textbf{The semantic metric completes the quality picture.} The
BERT-F1 ordering differs from the token-F1 ordering: Pure Hybrid is
semantically strongest ($0.865$) despite a mid-pack token-F1, Always-on is
second ($0.861$), and the Single-stage Router and the plain Oracle are
weakest ($0.735$, $0.771$). Combined evidence---mid EM, top BERT-F1---
indicates that hybrid evidence produces answers that are semantically
faithful but lexically divergent from the gold span, exactly the behaviour
token-F1 under-credits (Section~\ref{subsec:metrics}).

\noindent\textbf{Retrieval coverage now measures a genuine mechanism gap.}
After canonical identifiers are emitted at the retrieval stage, Hit@$1$ is
comparable across all systems: always-dense coverage remains the highest
($0.727$), the routed systems reach $0.508$--$0.588$, and the Single-stage
Router collapses to $0.215$---it often selects a defensible route yet
retrieves the wrong article, another symptom of operating without the
verifier. Closing the coverage gap of graph-mediated retrieval (e.g.\ by
cross-encoder re-ranking) is a concrete direction for future work.

\noindent\textbf{Pure Graph trades fluency for relational precision.} It
ties the best EM ($0.335$) but posts a lower F1 ($0.633$), reflecting
concise answers generated from highly structured graph context.

To locate \emph{where} each system earns its score,
Table~\ref{tab:per_class_f1} reports Answer~F1 separately by gold routing
class.

\begin{table}[!htbp]
\centering
\caption{Answer~F1 broken down by gold routing class on the strict benchmark
($300$ dense / $150$ graph / $150$ hybrid queries). Class-weighted means
reproduce the aggregate F1 of Table~\ref{tab:end_to_end}. Bold marks the best
value per column, excluding the two oracle rows}\label{tab:per_class_f1}
\begin{tabular}{lrrr}
\toprule
\textbf{System} & \textbf{Dense F1} & \textbf{Graph F1} & \textbf{Hybrid F1} \\
\midrule
Pure Vector      & 0.549 & \textbf{0.888} & \textbf{0.667} \\
Pure Graph       & \textbf{0.810} & 0.499 & 0.415 \\
Pure Hybrid      & 0.574 & 0.790 & 0.649 \\
Single-stage     & 0.560 & 0.377 & 0.331 \\
Two-stage (ours) & 0.646 & 0.509 & 0.607 \\
Always-on Stage~2 & 0.806 & 0.481 & 0.633 \\
\midrule
Oracle           & 0.466 & 0.515 & 0.649 \\
Oracle + Stage~2 & 0.661 & 0.499 & 0.640 \\
\bottomrule
\end{tabular}
\end{table}

The two fixed single-mechanism rows exhibit a striking and initially
counter-intuitive pattern: \emph{each pure mechanism scores highest on the
other mechanism's class}. Pure Graph attains the best F1 on dense-class
queries ($0.810$) while Pure Vector attains the best F1 on graph-class
queries ($0.888$); both rows are unchanged from the pre-fix run (their
outputs are deterministic and identical), so the inversion is a stable
property of the measurements, not a tabulation error. We interpret it as a
metric--prompt interaction rather than as evidence about mechanism
appropriateness: token-F1 is computed against a gold context, and the
citation-constrained generator quotes legal text verbatim, so a system
scores highly on a class whenever it happens to quote text with high lexical
overlap with that class's gold contexts---graph-class gold contexts are
article passages that dense retrieval reproduces nearly verbatim whenever
the target article is findable by embedding similarity, while concise
dense-class factoid questions penalise over-quotation. Exact match, which is
insensitive to quotation length, and BERT-F1, which is insensitive to
lexical form, both show the expected mechanism alignment. The consequence is
methodological: \emph{per-class token-F1 cannot adjudicate which retrieval
mechanism suits which query class}, which is why the semantic column of
Table~\ref{tab:end_to_end} is a required complement rather than an optional
extra (Section~\ref{subsec:metrics}). The adaptive rows nonetheless carry a
clear signal: the selective Two-stage router has the most even non-oracle
profile ($0.646/0.509/0.607$), and its gains over the Single-stage Router
concentrate exactly on the graph-class ($+0.132$) and hybrid-class
($+0.276$) queries where Stage~1 alone collapses ($0.377$, $0.331$), whereas
the always-on variant buys its aggregate lead mainly on the dense majority
class ($0.806$) at a slight cost on the graph class.

\subsection{RQ2: Routing Accuracy}\label{subsec:rq2}

A router is useful only if it selects the correct retrieval mechanism. We
evaluate routing from two complementary angles: the generalisation of the
deployed router on the paraphrase-hardened strict test split, and a
controlled comparison of classifier families under one identical
cross-validation protocol on the training data.

\paragraph{Generalisation on the strict split.}
Table~\ref{tab:routing_strict} reports the routing decision of the deployed
two-stage router on the $600$ strict test queries, whose surface templates
were paraphrased away from the training distribution. The router attains
$0.898$ overall accuracy and $0.900$ macro-F1 over the three retrieval
classes, with its weakest precision on the graph class ($0.790$), where
dense-class queries occasionally carry incidental relation cues. It
additionally emitted nine clarification requests ($1.5\%$) on answerable
queries---the only clarify false positives observed on this split.

\begin{table}[!htbp]
\centering
\caption{Deployed router on the strict $600$-query split: per-class recall
(R), precision (P), and F1 over the three retrieval classes. Overall accuracy
$0.898$; macro-F1 $0.900$. Nine clarify decisions ($1.5\%$) were emitted on
answerable queries}\label{tab:routing_strict}
\begin{tabular}{lrrrr}
\toprule
\textbf{Gold class} & \textbf{N} & \textbf{R} & \textbf{P} & \textbf{F1} \\
\midrule
Dense retrieval  & 300 & 0.890 & 0.964 & 0.926 \\
Graph traversal  & 150 & 0.853 & 0.790 & 0.820 \\
Hybrid reasoning & 150 & 0.960 & 0.947 & 0.953 \\
\midrule
\textbf{Overall / macro} & 600 & 0.901 & 0.900 & 0.900 \\
\bottomrule
\end{tabular}
\end{table}

Route~Acc in Table~\ref{tab:end_to_end} is lower ($0.623$ for the selective
system) because it scores the \emph{executed} route after Stage~2 overrides,
clarification decisions, and execution-time fallbacks, whereas
Table~\ref{tab:routing_strict} scores the routing \emph{decision}; the gap
between the two quantifies how often the verifier deliberately departs from
the gold taxonomy---departures that RQ1 shows to be quality-improving.

\paragraph{Classifier comparison under identical cross-validation.}
Table~\ref{tab:routing_cv} compares six classifiers under one stratified
five-fold protocol (identical folds, seed~$42$) over the $860$ training
samples ($800$ routing-labelled queries plus $60$ synthetic \texttt{clarify}
samples). Learned classifiers dominate the heuristic baselines by a wide
margin---at least $0.37$ accuracy and $0.44$ macro-F1 separate the weakest
learned model from the strongest rule-based one---confirming that routing
among evidence mechanisms cannot be reduced to a handful of surface rules.

\begin{table}[!htbp]
\centering
\caption{Routing classifiers under identical stratified five-fold
cross-validation on the training data (mean\,$\pm$\,std across folds).
Absolute values for the learned models are optimistic due to pre-split
paraphrase leakage (see text); the deployment-relevant figure is the
strict-split result of Table~\ref{tab:routing_strict}}\label{tab:routing_cv}
\begin{tabular}{lrr}
\toprule
\textbf{Classifier} & \textbf{CV Acc.} & \textbf{CV Macro-F1} \\
\midrule
Majority Route        & $0.4651\ \ci{0.0000}$ & $0.1587\ \ci{0.0000}$ \\
Keyword Rule Router   & $0.6081\ \ci{0.0182}$ & $0.3866\ \ci{0.0263}$ \\
Simple 3-Rule Router  & $0.6209\ \ci{0.0222}$ & $0.5422\ \ci{0.0308}$ \\
Logistic Regression   & $0.9953\ \ci{0.0044}$ & $0.9953\ \ci{0.0042}$ \\
Random Forest         & $0.9965\ \ci{0.0028}$ & $0.9966\ \ci{0.0029}$ \\
XGBoost (Stage~1, deployed) & $0.9942\ \ci{0.0037}$ & $0.9944\ \ci{0.0035}$ \\
\bottomrule
\end{tabular}
\end{table}

\paragraph{Leakage analysis: reconciling the routing figures.}
All three learned models exceed $0.99$ under cross-validation, which a
quantitative leakage analysis shows to be optimistic: because paraphrase
augmentation was applied \emph{before} splitting, $33.2\%$ of highly similar
query pairs (token Jaccard $\ge 0.7$) straddle the fold boundary, and the
mean Jaccard similarity of cross-fold pairs ($0.192$) is essentially equal to
that of within-fold pairs ($0.190$)---the folds do not isolate query
templates, so every learned model partially recognises test templates from
training. This resolves the apparent conflict between the routing figures
reported at different points of this project: the single-split and
cross-validated results ($\approx 0.99$) share the same template leakage and
should be read as upper bounds on separability given the features, whereas
the strict-split accuracy of $0.898$
(Table~\ref{tab:routing_strict})---measured on paraphrase-hardened queries
disjoint from training---is the honest generalisation estimate and the figure
we headline. Constructing template-grouped folds is left as an evaluation
improvement for future work.

\noindent\textbf{Feature reliance and efficiency.} The decision is driven
mainly by structural signals---the legal reference count, query length,
cross-document signals, and article reference count---rather than by spurious
lexical artefacts (Section~\ref{subsec:features};
Appendix~\ref{app:repro}, Table~\ref{tab:features}). The classifier forward
pass takes $\approx\!8.9$~ms on CPU versus $\approx\!53.3$~ms on GPU for a
PhoBERT classifier, making Stage~1 well suited to gating the expensive
Stage~2 verifier; the decisive evidence for the design, however, is the
end-to-end quality and latency trade-off of RQ1 and RQ5.

\subsection{RQ3: Ambiguity and Clarification Handling}\label{subsec:rq3}

Routing is safety-critical when the legal target is underspecified: retrieving
against an arbitrary document yields a confidently wrong answer. We evaluate
the deployed system on the $234$-query ambiguity benchmark ($156$ ambiguous
queries, $78$ answerable controls) and on the $160$-query multi-turn
conversation stress test ($8$ scenario types $\times$ $20$ dialogues).
Table~\ref{tab:ambiguity} summarises both benchmarks;
Table~\ref{tab:ambiguity_type} breaks the conversation results down by
scenario type.

\begin{table}[!htbp]
\centering
\caption{Clarification performance of the deployed router. P, R, and F1 are
computed over the \texttt{clarify} decision against the binary
ambiguous/answerable gold label; Overr.\ is the rate at which Stage~2
overrides the Stage~1 proposal}\label{tab:ambiguity}
\scriptsize
\setlength{\tabcolsep}{3pt}
\begin{tabular}{lrrrrrr}
\toprule
\textbf{Benchmark} & \textbf{Route Acc.} & \textbf{P} & \textbf{R} & \textbf{F1} & \textbf{S2 Trig.} & \textbf{S2 Overr.} \\
\midrule
Ambiguity ($N{=}234$)    & 0.761 & 0.945 & \textbf{1.000} & \textbf{0.972} & 85.0\% & 5.5\% \\
Conversation ($N{=}160$) & 0.725 & 0.857 & 0.960 & 0.906 & 91.9\% & 13.6\% \\
\bottomrule
\end{tabular}
\end{table}

\begin{table}[!htbp]
\centering
\caption{Conversation stress test by scenario type ($20$ dialogues each).
Clarify F1 is $0.000$ by construction for scenario types whose gold label is
a retrieval route rather than \texttt{clarify}}\label{tab:ambiguity_type}
\small
\begin{tabular}{lrrr}
\toprule
\textbf{Scenario type} & \textbf{Route Acc.} & \textbf{Clarify F1} & \textbf{S2 Trig.} \\
\midrule
Clarify, no history          & 1.000 & 1.000 & 1.00 \\
Missing entity               & 1.000 & 1.000 & 1.00 \\
Multi-interpretation         & 1.000 & 1.000 & 1.00 \\
Irrelevant history           & 0.950 & 0.974 & 0.95 \\
Conflicting history          & 0.850 & 0.919 & 0.90 \\
Answerable with history      & 0.500 & 0.000 & 0.85 \\
Clear graph/hybrid (control) & 0.400 & 0.000 & 0.75 \\
Clear dense (control)        & 0.100 & 0.000 & 0.90 \\
\bottomrule
\end{tabular}
\end{table}

We make the following observations.

\noindent\textbf{Every ambiguous query is intercepted.} On the ambiguity
benchmark the router attains clarify recall $1.000$ at precision $0.945$
(F1~$=0.972$): none of the $156$ ambiguous queries reaches retrieval, at the
cost of spuriously clarifying roughly nine of the $78$ answerable controls
($11.5\%$). For a legal assistant this operating point is the safer side of
the trade---a missed ambiguity produces a confidently wrong answer grounded
in the wrong document, whereas a spurious clarification costs a single
dialogue turn. The contrast between the $11.5\%$ false-positive rate on the
diagnostic's controls and the $1.5\%$ clarify rate on the ambiguity-free
strict split (Section~\ref{subsec:rq2}) shows, however, that precision
degrades as the ambiguity base rate of the surrounding traffic rises.

\noindent\textbf{Semantic underspecification is detected, not just surface
cues.} The two scenario types that carry no surface marker of
ambiguity---missing entity and multi-interpretation, where the query is
grammatically complete but legally underdetermined---are detected perfectly
(clarify F1~$=1.000$, trigger rate $100\%$). The indicator families targeting
generic legal needs without a concrete anchor ($\psi_6$) and broad
subject--predicate combinations without a domain anchor ($\psi_8$), combined
with the weakly supervised Stage~1 clarify class, therefore capture semantic
underspecification rather than merely pronoun surface cues.

\noindent\textbf{History conditions are handled asymmetrically.} Dialogues
whose history is irrelevant or conflicting are clarified near-perfectly
(F1~$=0.974$ and $0.919$), and clarification without history is perfect. The
weak case is \emph{answerable with history}: only half of the queries whose
referent the resolver can recover are subsequently routed to the correct
retrieval mechanism, the remainder being either spuriously clarified or
dispatched to the wrong route.

\noindent\textbf{Clear queries embedded in conversation are misrouted.} The
two control scenario types expose the principal remaining weakness:
self-contained dense queries inside a dialogue are routed correctly only
$10\%$ of the time (graph/hybrid controls: $40\%$), even though Stage~2 is
triggered on $75$--$90\%$ of them. Conversational context evidently inflates
the ambiguity and trigger signals of queries that need no context at all,
dragging them off their correct route; overall conversation route accuracy
($0.725$) is therefore bounded by control misrouting, not by clarify
detection (F1~$=0.906$). Suppressing conversational features whenever the
resolver returns \texttt{not\_needed}---that is, when the query is
self-contained---is the natural mitigation and is left to future work.

\noindent\textbf{The verifier mostly confirms Stage~1.} Stage~2 overrides its
Stage~1 proposal on only $5.5\%$ of ambiguity-benchmark queries and $13.6\%$
of conversation queries, concentrating its interventions where history must
be adjudicated. As a concrete example, given history mentioning
Decree ``100/2019/ND-CP'', the follow-up ``Is that document still in
effect?'' is resolved and routed to graph traversal, whereas with empty
history the same query is routed to \texttt{clarify}.

\subsection{RQ4: Threshold Ablation}\label{subsec:rq4}

The router exposes two key thresholds: the confidence threshold
$\tau_{\mathrm{conf}}$ that gates Stage~2, and the clarification override
threshold $\tau_{\mathrm{clar}}$. We study their effect on the development set
in Table~\ref{tab:threshold_ablation}.

\begin{table}[!htbp]
\centering
\caption{Threshold ablation on the development set}\label{tab:threshold_ablation}
\begin{tabular}{lrrr}
\toprule
\textbf{Parameter} & \textbf{Value} & \textbf{Acc.} & \textbf{S2 rate} \\
\midrule
$\tau_{\mathrm{conf}}$ & 0.50 & 0.5200 & 0.8600 \\
$\tau_{\mathrm{conf}}$ & 0.70 & 0.5200 & 0.8600 \\
$\tau_{\mathrm{conf}}$ & 0.95 & 0.5200 & 0.8600 \\
$\tau_{\mathrm{conf}}$ & 0.99 & 0.5200 & 0.8600 \\
\midrule
$\tau_{\mathrm{clar}}$ & 0.50 & 0.1400 & 0.8600 \\
$\tau_{\mathrm{clar}}$ & 0.70 & 0.5200 & 0.8600 \\
$\tau_{\mathrm{clar}}$ & 0.95 & 0.6700 & 0.8600 \\
$\tau_{\mathrm{clar}}$ & 0.99 & 0.6700 & 0.8600 \\
\bottomrule
\end{tabular}
\end{table}

We make the following observations.

\noindent\textbf{$\tau_{\mathrm{conf}}$ is inert on this development set.}
Sweeping $\tau_{\mathrm{conf}}$ from $0.50$ to $0.99$ changes neither
accuracy ($0.5200$) nor the Stage~2 trigger rate ($0.8600$): the trigger is
saturated by the confidence-independent conditions---the ambiguity forcing
rule ($\phi_{\mathrm{amb}}\ge\tau_{\mathrm{force}}$), the multi-hop
condition, and the mandatory adjudication of Stage~1 clarify
proposals---which already fire on $86\%$ of development queries. Cost
control on ambiguity-heavy traffic must therefore come from those
conditions, not from the confidence gate alone; on the cleaner strict split
the effective trigger rate is lower ($79.2\%$), indicating that the
confidence gate matters only when the other conditions are quiet.

\noindent\textbf{$\tau_{\mathrm{clar}}$ has a strong, monotone accuracy
effect.} A permissive override of $0.50$ collapses development accuracy to
$0.1400$ by over-clarifying answerable queries; tightening to $0.70$ recovers
$0.5200$, and $0.95$ plateaus at $0.6700$ (unchanged at $0.99$). The default
$\tau_{\mathrm{clar}}=0.80$ sits on the rising part of this curve,
reflecting the opposing pull of the two objectives: the value that maximises
routing accuracy on answerable queries is not the one that maximises clarify
recall on the ambiguity benchmark, so the two must be jointly calibrated
rather than tuned in isolation.

To connect this ablation to the RQ5 cost argument,
Table~\ref{tab:trigger_policy} contrasts the three verification
policies---always-off, selective (ours), and always-on---on the strict
benchmark. The three policies now separate cleanly: always-off is cheapest
and weakest, always-on is the quality endpoint, and the selective policy
recovers $471$~ms ($4.6\%$) of latency at a statistically significant cost
of $0.080$~F1. On this deployment, where a local verifier call is cheap
relative to generation, always-on is the preferable operating point; the
selective policy is the right choice when verifier invocations carry a
dominant marginal cost (remote APIs, per-token pricing, or higher traffic),
as quantified by the $119\times$ routing-stage cost ratio of RQ5.

\begin{table}[!htbp]
\centering
\caption{Verification-policy comparison on the strict
benchmark}\label{tab:trigger_policy}
\begin{tabular}{lrrr}
\toprule
\textbf{Policy} & \textbf{S2 rate} & \textbf{F1} & \textbf{Avg lat (ms)} \\
\midrule
Always-off (Single-stage) & 0.0\%   & 0.457 & 7{,}262 \\
Selective (Two-stage, ours) & 79.2\% & 0.602 & 9{,}799 \\
Always-on & 100.0\% & \textbf{0.682} & 10{,}270 \\
\bottomrule
\end{tabular}
\end{table}

\subsection{RQ5: Cost--Quality Trade-off}\label{subsec:rq5}

Stage~2 is powerful but expensive, so the central design question is how much
verification to buy. Table~\ref{tab:cost} reports the latency--quality
analysis for the post-fix run.

\begin{table}[!htbp]
\centering
\caption{Cost--quality trade-off on the post-fix run. S2 is the Stage~2
trigger rate; the Oracle latency reflects gold-label routing without Stage~2
overhead}\label{tab:cost}
\small
\begin{tabular}{lrrrrr}
\toprule
\textbf{System} & \textbf{F1} & \textbf{Mean (ms)} & \textbf{p50 (ms)} & \textbf{p95 (ms)} & \textbf{S2} \\
\midrule
Pure Vector  & 0.663 & 10{,}914 & 11{,}542 & 33{,}776 & -- \\
Pure Graph   & 0.633 & \textbf{6{,}034} & \textbf{4{,}206} & \textbf{19{,}383} & -- \\
Pure Hybrid  & 0.647 & 7{,}375 & 5{,}719 & 20{,}511 & -- \\
Single-stage & 0.457 & 7{,}262 & 6{,}117 & 19{,}896 & 0.0\% \\
Two-stage (ours) & 0.602 & 9{,}799 & 8{,}159 & 21{,}758 & 79.2\% \\
Always-on Stage~2 & \textbf{0.682} & 10{,}270 & 8{,}950 & 20{,}881 & 100.0\% \\
\midrule
Oracle       & 0.524 & 7{,}548 & 5{,}851 & 19{,}505 & -- \\
Oracle + Stage~2 & 0.615 & 9{,}920 & 8{,}325 & 20{,}529 & 79.2\% \\
\bottomrule
\end{tabular}
\end{table}

We make the following observations.

\noindent\textbf{Always-off verification is cheap but weak.} The Single-stage
Router skips Stage~2 entirely and runs at $7{,}262$~ms on average, but pays
the largest quality penalty (F1~$=0.457$); Stage~2 is the sole source of the
$+0.145$ (selective) and $+0.225$ (always-on) F1 gaps.

\noindent\textbf{At the routing stage, Stage~2 is two orders of magnitude more
expensive than Stage~1.} A routing-stage microbenchmark, measured separately
from the end-to-end runs, puts one Stage~2 verifier call at
$\approx\!4{,}017$~ms against $\approx\!34$~ms for the full Stage~1
path---roughly a $119\times$ ratio. When the verifier is served remotely or
priced per token, this ratio makes selective triggering the principled
design; on the present local deployment the ratio is absorbed by the
generation-dominated budget analysed next.

\noindent\textbf{End-to-end, verification costs seconds and buys tenths of
F1.} Selective verification adds $2{,}537$~ms ($+35\%$) over the
Single-stage Router for $+0.145$~F1; going from selective to always-on adds
only a further $471$~ms ($+4.6\%$) for another $+0.080$~F1, because $79.2\%$
of queries already pay the verifier. The marginal quality per millisecond
therefore \emph{improves} at the always-on end of the dial, which is why RQ4
recommends always-on for this deployment.

\noindent\textbf{Latency is dominated by generation, not by routing.}
Table~\ref{tab:latency_decomp} decomposes the Two-stage average: Stage~1
contributes $\approx\!34$~ms and the amortised verifier
$\approx\!3{,}181$~ms ($4{,}017$~ms per triggered call at a $79.2\%$ rate),
against $\approx\!2{,}500$~ms of retrieval and $\approx\!4{,}084$~ms of
answer generation. The routing-efficiency argument for Stage~1 is therefore
that routing adds negligible overhead, not that it dominates latency.

\begin{table}[!htbp]
\centering
\caption{Latency decomposition (ms) for the Two-stage system, separating
routing overhead from retrieval and generation}\label{tab:latency_decomp}
\begin{tabular}{lr}
\toprule
\textbf{Component} & \textbf{Avg (ms)} \\
\midrule
Stage~1 routing            & $\approx 34$ \\
Stage~2 verification (amortised at $79.2\%$; $\approx 4{,}017$ per call) & $\approx 3{,}181$ \\
Retrieval (dense / graph / hybrid) & $\approx 2{,}500$ \\
Answer generation          & $\approx 4{,}084$ \\
\midrule
End-to-end (measured)      & 9{,}799 \\
\bottomrule
\end{tabular}
\end{table}

\noindent\textbf{The trigger conditions, not the confidence gate, are the
cost dial.} Because $\tau_{\mathrm{conf}}$ is saturated by the ambiguity and
relational conditions (Section~\ref{subsec:rq4}), tuning the trigger rate on
ambiguity-heavy traffic requires adjusting those conditions
($\tau_{\mathrm{force}}$, the multi-hop rule, and the clarify-proposal
policy); the always-off and always-on endpoints in
Table~\ref{tab:trigger_policy} bracket the achievable operating points.

\subsection{Overall Discussion}\label{subsec:discussion}

The experiments support the central thesis that Vietnamese legal QA should
not use a single fixed retrieval strategy, and the post-fix run sharpens
where the benefit lies. First, the cleanest quantitative result is causal:
holding routing at gold labels and toggling only the verifier isolates a
$+0.091$ token-F1 and $+0.052$ BERT-F1 contribution of Stage~2 reasoning
with non-overlapping confidence intervals, confirming that the verifier
functions as a chain-of-thought step during answer generation rather than as
a route corrector---indeed, executed-route agreement falls monotonically
($0.737\!\to\!0.623\!\to\!0.455$) while quality rises, peaking exactly where
agreement is lowest. Second, the framework's quality-optimal operating point
is always-on verification: F1~$=0.682$, statistically indistinguishable from
and nominally above the strongest fixed baseline ($0.663$), tied for the
best exact match ($0.335$, $2.2\times$ always-dense), and within $0.004$ of
the best semantic score---no fixed strategy matches this profile, and only
the routed systems can refuse to guess. The selective policy sits
significantly below that point on token-F1 ($0.602$) while saving only
$4.6\%$ latency here; its case rests on deployments where verifier cost
dominates, per the $119\times$ routing-stage ratio. Third, the per-class
analysis exposes a metric--prompt interaction---each pure mechanism scores
highest on the other mechanism's class under token-F1---which cautions
against reading cross-system token-F1 differences as mechanism quality; the
semantic column shows a different ordering (hybrid evidence strongest) and
the causal ablation replicates on it. Fourth, the clarification analysis
remains the cleanest behavioural result: ambiguous queries---including the
semantically underspecified missing-entity and multi-interpretation
types---are intercepted with perfect recall, while the misrouting of clear
queries embedded in conversation is now the principal open weakness,
motivating context-gated suppression of conversational features for
self-contained queries.

%% ===== limitations =====================================================
\section{Limitations and Future Work}\label{sec:limitations}
% ===============================================================

Several limitations qualify the results above and define the agenda for future
work.

\noindent\textbf{Synthetic, metadata-derived benchmark.} Route labels are
derived from hop-count and cross-document structure rather than annotated by
legal experts, and the benchmark---including the ambiguity and conversation
stress tests---is machine-generated. Although paraphrasing mitigated
template-based label leakage on the strict split, the distribution remains
synthetic and should be validated against a human-authored set before strong
external-validity claims are made.

\noindent\textbf{Weakly supervised clarify class.} The strict split contains
no gold \texttt{clarify} label, so the Stage~1 clarify class is trained on
only $60$ synthetic templates and is deliberately treated as a proposal that
Stage~2 must confirm. End-to-end clarify routing therefore rests on the
verifier and the ambiguity detector rather than on strongly supervised
classification.

\noindent\textbf{In-conversation routing of self-contained queries.} On the
conversation stress test, clearly answerable control queries embedded in a
dialogue are routed correctly only $10\%$ (dense) and $40\%$ (graph/hybrid)
of the time: conversational context inflates the ambiguity and trigger
signals of queries that need no context, and the resolver's
\texttt{not\_needed} status does not yet suppress those signals. Conversation
route accuracy ($0.725$) is bounded by this misrouting rather than by
clarify detection, and the clarify false-positive rate rises from $1.5\%$ on
the ambiguity-free strict split to $11.5\%$ on the diagnostic's answerable
controls as the ambiguity base rate of the surrounding traffic increases.

\noindent\textbf{Sparse legal-effect relations and explicit-only hybrids.}
Cross-document legal-effect relations (\textsc{AMENDS}, \textsc{REPLACES})
remain comparatively sparse in the graph, and the benchmark covers only
\emph{explicit} hybrid questions that name the documents to combine, not
\emph{implicit} ones that require discovering the link first.

\noindent\textbf{Retrieval coverage of graph-mediated evidence.} With
canonical \texttt{law::document::article} identifiers now emitted at the
retrieval stage and matched strictly at evaluation, Hit@$1$ is comparable
across all systems and reveals a genuine coverage gap: the routed
configurations reach $0.508$--$0.588$ against $0.727$ for always-dense
retrieval, and the Single-stage Router falls to $0.215$. Graph and hybrid
evidence paths therefore locate the gold article less reliably than dense
search even when the route is appropriate; cross-encoder re-ranking of graph
candidates is the natural remedy. Mean reciprocal rank, now computable from
the full ranking logs, remains to be extracted.

\noindent\textbf{Aggregate token-F1 depends on the trigger policy.} On the
dense-dominated benchmark the selective router ($0.602$) is
\emph{significantly below} the Pure Vector baseline ($0.663$) on aggregate
token-F1, with non-overlapping bootstrap intervals; statistical parity with
(and nominal superiority over) the strongest fixed baseline is achieved only
at the always-on operating point ($0.682$). The framework's contribution is
therefore argued on the causal Stage~2 ablation, on exact match
($2.2\times$ always-dense at the always-on point), on the semantic metric,
and on clarification safety---and any claim about the selective policy must
carry its measured quality cost.

\noindent\textbf{Metric validity.} Token-F1 against a gold context rewards
lexical overlap, and the citation-constrained generator quotes legal text
verbatim; the combination produces the class-wise inversion of
Table~\ref{tab:per_class_f1}, in which each pure mechanism scores highest on
the \emph{other} mechanism's query class. Cross-system token-F1 differences
therefore partly measure quotation overlap rather than answer correctness,
and per-class token-F1 cannot adjudicate mechanism appropriateness. The
aggregate BERT-F1 column mitigates this at the system level---its ordering
differs from token-F1's, and the causal ablation replicates on it---but
per-class semantic scores have not been extracted, so the class-level
analysis still rests on the biased lexical metric.

\noindent\textbf{Scope of the chain-of-thought claim.} The
\textit{Oracle~+~Stage~2} ablation isolates a causal $+0.091$ token-F1 and
$+0.052$ BERT-F1 contribution of Stage~2 reasoning under fixed gold routing,
but the estimate derives from a single generator (\texttt{Qwen3-32B-AWQ}), a
single prompt family, and deterministic decoding on one benchmark; its
magnitude may not transfer across models, prompts, or legal domains.

\noindent\textbf{Statistical reporting gaps.} Headline F1 comparisons carry
bootstrap $95\%$ confidence intervals ($B=1000$), but no paired significance
tests are reported. We report retrieval coverage through Hit@$1$ rather than
mean reciprocal rank: an auxiliary MRR computation was consistent with Hit@$1$
for the router-mediated systems but produced values below Hit@$1$ for two
fixed baselines---a violation of the identity $\mathrm{MRR}\ge\mathrm{Hit@1}$
that traces to a separate rank-logging path for those baselines---so we omit
MRR pending a unified ranking log rather than report an inconsistent column.

\noindent\textbf{Routing-accuracy evaluation under template leakage.} Because
paraphrase augmentation preceded splitting, $33.2\%$ of highly similar query
pairs cross the cross-validation fold boundary, inflating the absolute CV
scores of all learned routers to $\approx 0.99$. We therefore headline the
paraphrase-hardened strict-split accuracy ($0.898$) instead, and leave
template-grouped fold construction as an evaluation improvement for future
work.

\noindent\textbf{No empirical comparison to complexity-based adaptive
retrieval.} Adaptive-RAG~\cite{jeong2024adaptive} is the closest prior method
but is not compared empirically, because it routes among reasoning depths for
general-domain English QA rather than among evidence mechanisms for Vietnamese
legal QA. A faithful re-implementation of its complexity router as an
additional baseline---or an explicit statement of why the settings are not
directly comparable---would strengthen the positioning.

\noindent\textbf{Efficiency and hardware dependence.} A Stage~2 verifier call
is roughly $119\times$ the cost of the full Stage~1 routing path in the
routing-stage microbenchmark, and all latencies depend on local vLLM
hardware; end-to-end latency is dominated by generation rather than by
routing, so the efficiency argument is about negligible routing overhead
rather than end-to-end speed.

Future directions include: a three-level cascade
XGBoost~$\rightarrow$~PhoBERT~$\rightarrow$~LLM to preserve routing accuracy
under the linguistic diversity of naturally phrased queries; cross-encoder
re-ranking of graph and hybrid candidates to close the retrieval-coverage
gap to dense search;
integration of manually reviewed \texttt{clarify} samples and implicit hybrid
questions into the training set; context-gated suppression of conversational
features whenever the history resolver marks a query as self-contained;
stronger legal-relation extraction for
\textsc{AMENDS} and \textsc{REPLACES}; a complexity-router baseline in the
spirit of Adaptive-RAG; semantic and human evaluation of legal answer
correctness and citation faithfulness; and template-grouped cross-validation
for leakage-free router comparison.

%% ===== conclusion ======================================================
\section{Conclusion}\label{sec:conclusion}
% ===============================================================

We presented a reasoning-aware adaptive router for Vietnamese legal Hybrid
GraphRAG that combines a dense vector index of $68{,}835$ legal text chunks, a
Neo4j legal knowledge graph of $804{,}322$ nodes and $1{,}968{,}621$ edges
built from the national legal codification, an XGBoost Stage~1 classifier, and
a selective LLM Stage~2 verifier, dispatching each query to dense retrieval,
graph traversal, hybrid reasoning, or clarification according to its reasoning
demand and ambiguity level.

The central quantitative findings on the strict $600$-query benchmark, all
reported with bootstrap $95\%$ confidence intervals, are as follows. A
controlled \textit{Oracle~+~Stage~2} ablation that holds routing at gold
labels while enabling the verifier isolates a causal $+0.091$ token-F1 and
$+0.052$ BERT-F1 contribution of Stage~2 reasoning ($0.615$ vs $0.524$
token-F1), confirming that the second stage acts as a chain-of-thought step
during answer generation rather than as a route corrector---executed-route
agreement with the gold labels falls monotonically as the verifier
intervenes more ($0.737\!\to\!0.623\!\to\!0.455$) while quality rises at
every step. At its always-on operating point the framework reaches
F1~$=0.682$, statistically indistinguishable from and nominally above the
strongest fixed baseline (Pure Vector, $0.663$), tied for the best exact
match ($0.335$, $2.2\times$ always-dense), and within $0.004$ of the best
semantic BERT-F1; no fixed strategy matches this profile across all three
answer metrics. The selective policy at a $79.2\%$ trigger rate trades
$0.080$~F1 for a $4.6\%$ latency saving, exceeding the Single-stage Router
by $+0.145$~F1 and the plain Oracle by $+0.078$. The Stage~1 classifier
generalises to the paraphrase-hardened strict split at $0.898$ accuracy and
$0.900$ macro-F1 at millisecond latency, far above rule-based routers under
an identical cross-validation protocol.

On the ambiguity benchmarks, the deployed router intercepts every ambiguous
query before retrieval---clarify recall $1.000$ at precision $0.945$
(F1~$=0.972$) on the $234$-query diagnostic---and sustains clarify
F1~$=0.906$ on the $160$-query multi-turn stress test, detecting even the
semantically underspecified missing-entity and multi-interpretation scenarios
perfectly. The remaining weakness is the routing of clearly answerable
queries embedded in conversation, where contextual trigger signals pull
self-contained controls off their correct route.

Closing the retrieval-coverage gap of graph-mediated evidence through
cross-encoder re-ranking, suppressing conversational features for
self-contained queries, implementing a complexity-router baseline, improving
routing robustness through a three-level cascade, and validating the
benchmarks with human-reviewed legal queries remain the primary directions
for future investigation.


% Back matter, declarations, appendix
%% ===== declarations ====================================================
\backmatter

\bmhead{Acknowledgements}
The authors thank the supervising faculty at the School of Information and
Communication Technology, Hanoi University of Science and Technology, for their
guidance throughout this work.

\section*{Declarations}


\noindent\textbf{Funding.} This research received no specific grant from any
funding agency in the public, commercial, or not-for-profit sectors.

\noindent\textbf{Competing interests.} The authors declare that they have no
competing interests.

\noindent\textbf{Ethics approval and consent to participate.} Not applicable.
This study does not involve human participants, their identifiable data, or
biological material. The corpus consists solely of publicly available
Vietnamese legal documents from the official electronic legal codification.

\noindent\textbf{Consent for publication.} Not applicable.

\noindent\textbf{Data availability.} The Vietnamese legal QA routing dataset,
the ambiguity and conversation benchmarks, and the derived routing labels are
publicly available at \url{https://REPLACE-WITH-YOUR-DATA-URL}.

\noindent\textbf{Code availability.} The routing pipeline, feature extractor,
Stage~1 training scripts, Stage~2 verifier prompts, and evaluation harness are
publicly available at \url{https://REPLACE-WITH-YOUR-CODE-URL}.

\noindent\textbf{Author contributions.} N.D.T.\ designed and implemented the
system, conducted all experiments, and drafted the manuscript. P.V.H.\ and
C.T.\ supervised the research, reviewed the experimental design, and revised
the manuscript. All authors read and approved the final version.

%% ===== appendix ========================================================
\appendix

\section{Reproducibility Details}\label{app:repro}

This appendix documents the exact feature set, the ambiguity indicators and
their weights, the deterministic history-resolution rules, the evidence
identifier normalisation used for retrieval metrics, the generation and
clarification prompts, and the Stage~1 training configuration, so that the
router can be reproduced from the corpus alone.

\subsection{Stage~1 Feature Set}\label{app:features}

Table~\ref{tab:features} enumerates the $n_f=27$ features consumed by the
Stage~1 classifier of Eq.~\eqref{eq:stage1}, in the canonical vector order
shared by training and inference: features $0$--$15$ are lexical surface
features from the Vietnamese legal feature extractor, and features
$16$--$26$ are reasoning-aware metrics from the query-complexity analyser.
Cue lists are abbreviated; the four count features that drive routing use the
following patterns, verified against the feature-extractor source.
\texttt{graph\_keyword\_count} matches multi-hop relation cues glossed as
``related to'', ``between'', ``affects'', ``impacts'', ``relation'',
``leads to'', and ``compare''. \texttt{legal\_reference\_count} matches the
numbered document form \texttt{NN/YYYY/CODE} together with the document-type
nouns glossed as ``Decree'', ``Circular'', ``Law'', and ``Code''.
\texttt{dieu\_reference\_count} matches article and clause references
(``Article~N'', ``Clause~N''). \texttt{cross\_doc\_signal\_count} matches
cross-document cues glossed as ``the documents'', ``multiple laws'',
``simultaneously'', ``combine'', and ``both''.

\begin{table}[!htbp]
\centering
\caption{Stage~1 feature vector ($n_f=27$) in canonical order. Type ``bin''
denotes a binary indicator. Cue bases are given as English glosses of the
Vietnamese surface forms the extractor matches against}\label{tab:features}
\footnotesize
\begin{tabularx}{\columnwidth}{r l l Y}
\toprule
\# & \textbf{Feature} & \textbf{Type} & \textbf{Definition / cue basis} \\
\midrule
\multicolumn{4}{l}{\emph{Lexical features (Vietnamese legal feature extractor)}}\\
0 & \texttt{query\_length\_chars} & int & character count of the query \\
1 & \texttt{query\_length\_words} & int & whitespace token count \\
2 & \texttt{graph\_keyword\_count} & int & multi-hop relation cues (``related to'', ``between'', ``affects'', ``impacts'', ``leads to'', ``compare'') \\
3 & \texttt{legal\_reference\_count} & int & document references: numbered form \textit{NN/YYYY/CODE} and type nouns (``Decree'', ``Circular'', ``Law'', ``Code'') \\
4 & \texttt{dieu\_reference\_count} & int & article/clause references (``Article~N'', ``Clause~N'') \\
5 & \texttt{multi\_article\_ref} & bin & $1$ iff \texttt{dieu\_reference\_count} $\ge 2$ \\
6 & \texttt{is\_yes\_no\_question} & bin & yes/no interrogative endings (``\ldots or not'', ``is it true'', ``right?'') \\
7 & \texttt{is\_factoid\_question} & bin & factoid interrogatives (``when'', ``who'', ``how much'', ``where'') \\
8 & \texttt{cross\_doc\_signal\_count} & int & cross-document cues (``the documents'', ``multiple laws'', ``simultaneously'', ``combine'', ``both'') \\
9 & \texttt{has\_conditional\_structure} & bin & conditionals (``if'', ``in case'', ``provided that'', ``unless'') \\
10 & \texttt{has\_negation} & bin & negation (``not'', ``not yet'', ``no'', ``prohibit'') \\
11 & \texttt{has\_comparison} & bin & comparison/contrast (``compare'', ``differ'', ``same'', ``more'', ``less'') \\
12 & \texttt{has\_enumeration} & bin & enumeration markers (``include'', ``comprise'', ``list'') \\
13 & \texttt{has\_procedure\_marker} & bin & procedure markers (``procedure'', ``dossier'', ``process'', ``submit'', ``register'') \\
14 & \texttt{ambiguity\_score} & float & $\phi_{\mathrm{amb}}$ from the ambiguity detector, Eq.~\eqref{eq:ambiguity_score} \\
15 & \texttt{has\_multi\_hop\_connector} & bin & temporal/causal connectors (``after'', ``before'', ``leads to'', ``therefore'') \\
\midrule
\multicolumn{4}{l}{\emph{Reasoning features (query-complexity analyser)}}\\
16 & \texttt{complexity\_level} & int & overall complexity, $1$ (simple) to $3$ (complex) \\
17 & \texttt{sub\_question\_count} & int & decomposable sub-questions (clause and conjunction breaks) \\
18 & \texttt{law\_specificity} & float & specificity of law references, $0$ (vague) to $1$ (fully specified) \\
19 & \texttt{conditional\_depth} & int & nesting depth of conditional clauses \\
20 & \texttt{multi\_hop\_verb\_count} & int & verbs signalling multi-hop traversal \\
21 & \texttt{comparative\_depth} & int & depth/count of comparison operations \\
22 & \texttt{authority\_chain\_count} & int & authority-delegation patterns (``assign to'', ``authorise'', ``delegate'') \\
23 & \texttt{legal\_effect\_count} & int & legal-effect terms (``effect'', ``repeal'', ``amend'', ``replace'') \\
24 & \texttt{procedural\_count} & int & procedural step markers (``step'', ``stage'', ``approve'') \\
25 & \texttt{multi\_entity\_relation\_count} & int & multi-entity relation expressions \\
26 & \texttt{entity\_count} & int & legal entity mentions (legal-entity pattern) \\
\bottomrule
\end{tabularx}
\end{table}

\subsection{Ambiguity Indicators and Weights}\label{app:ambiguity}

Table~\ref{tab:ambiguity_features} lists the $m=11$ indicators
$\psi_i(q_t,H,r_H)$ and weights $w_i$ that define the ambiguity score of
Eq.~\eqref{eq:ambiguity_score},
\[
\phi_{\mathrm{amb}}=0 \ \text{if}\ r_H=\texttt{resolved};\qquad
\phi_{\mathrm{amb}}=\max_i\, w_i\,\psi_i \ \text{otherwise},
\]
so that a single strong indicator suffices to cross the Stage~2 trigger
threshold $\tau_{\mathrm{force}}=0.60$ while weak indicators cannot. Two
indicator families are \emph{graded} rather than binary: vague legal
references contribute $\min(1.0,\,0.30\cdot n_{\mathrm{vague}})$ (rows
$\psi_4$/$\psi_5$ show the one- and two-match values), and same-type entity
conflicts contribute $\min(0.5,\,0.15\cdot n_{\mathrm{conf}})$. The values
below are extracted directly from the detector's weight table.

\begin{table}[!htbp]
\centering
\caption{Ambiguity indicators $\psi_i$ and weights $w_i$ (extracted from the
ambiguity-detector source)}\label{tab:ambiguity_features}
\footnotesize
\begin{tabularx}{\columnwidth}{r l Y r}
\toprule
$i$ & \textbf{Indicator} & \textbf{Condition} & $w_i$ \\
\midrule
1 & \texttt{pronoun\_no\_history} & unresolved pronoun/demonstrative, empty history & 0.90 \\
2 & \texttt{pronoun\_bad\_history} & pronoun with conflicting or irrelevant history & 0.85 \\
3 & \texttt{contextual\_reference} & contextual reference flagged by the history resolver & 0.80 \\
4 & \texttt{vague\_reference\_single} & one vague legal reference (e.g.\ ``regulation'', ``that law'') & 0.30 \\
5 & \texttt{vague\_reference\_multi} & two or more vague legal references (graded) & 0.60 \\
6 & \texttt{missing\_entity\_no\_intent} & generic legal need with no concrete legal anchor & 0.82 \\
7 & \texttt{missing\_entity\_with\_intent} & generic need but domain context present (weak) & 0.10 \\
8 & \texttt{multi\_interpretation} & broad subject and broad predicate, no domain anchor & 0.85 \\
9 & \texttt{short\_no\_entity} & query under five words with no legal entity & 0.70 \\
10 & \texttt{incomplete\_context} & contextual follow-up with no usable history & 0.90 \\
11 & \texttt{entity\_conflict} & two or more entities of the same type (graded) & 0.50 \\
\bottomrule
\end{tabularx}
\end{table}

\subsection{History-Resolution Rules}\label{app:resolver}

The deterministic resolver $\mathrm{Resolve}(q_t,H)$ of
Eq.~\eqref{eq:history_resolver} is intentionally conservative: any outcome
other than \texttt{resolved} or \texttt{not\_needed} forces Stage~2, so the
resolver never silently binds an uncertain referent. It assigns the status
$r_H$ by the following ordered rules, returning the first that matches:
\begin{enumerate}
\item If $q_t$ contains no contextual-reference pattern (no demonstrative,
      pronoun, or elliptical reference), return \texttt{not\_needed}: the
      query is self-contained and history is not consulted.
\item Otherwise, if $H=\emptyset$, return \texttt{no\_history}.
\item Otherwise, if $H$ yields no legal referent whose type matches the
      reference in $q_t$, return \texttt{irrelevant\_history}.
\item Otherwise, if two or more candidate referents of the same type carry
      comparably high confidence while $q_t$ uses a singular demonstrative,
      return \texttt{conflicting\_history}.
\item Otherwise, bind the single highest-confidence referent and return
      \texttt{resolved}.
\end{enumerate}
The \texttt{no\_history}, \texttt{irrelevant\_history}, and
\texttt{conflicting\_history} outcomes each raise the ambiguity signal and
route through Stage~2 to clarification, whereas \texttt{resolved} unblocks the
retrieval route implied by the bound referent.

\subsection{Evidence Identifier Normalisation}\label{app:idnorm}

Retrieval metrics compare identifiers from four historically incompatible
schemes: national document numbers (e.g.\ \texttt{77/2026/ND-CP}), short
decision and consolidated-document codes (\texttt{QD-}, \texttt{VBHN-}
prefixes), composite \texttt{law::article} strings, and codification
structural codes (e.g.\ \texttt{19.2.TT.10.15}). A deterministic normaliser
maps every retrieved and gold reference to the canonical form
\texttt{law::<document\_code>::article::<number>} after Unicode NFC
normalisation and homoglyph repair, extracting the document code and the
article number independently. Matching supports three modes: \emph{strict}
(document code and article number must both agree; used for all reported
Hit@$k$), \emph{article-only}, and \emph{document-only}. References the
normaliser cannot resolve to the gold scheme are counted as misses and
flagged as unresolvable rather than fuzzily matched. During development this
flagging exposed that the hybrid evidence handler logged document codes
without article numbers---unlike the dense path, which extracted
both---collapsing strict-mode Hit@$1$ for every configuration touching the
hybrid path; porting the article-extraction step into the hybrid handler
restored comparable coverage measurement across all systems
(Table~\ref{tab:end_to_end}).

\subsection{Generation, Clarification, and Verifier Prompts}\label{app:prompt}

\paragraph{Answer-generation system prompt.}
Answers are generated under the following citation-constrained system prompt
(translated into English), which scopes verbatim quotation to the clause or
point directly answering the question, permits a short direct answer before
the supporting quotation, and fixes a single not-found sentence.

\begin{verbatim}
You are a legal assistant specialising in Vietnamese
law. Your task is to answer the question by QUOTING
VERBATIM the directly relevant statutory provision in
[LEGAL DOCUMENTS]. Do NOT paraphrase and do NOT
summarise the legal content in your own words.

MANDATORY RULES:
1. Quote ONLY verbatim from [LEGAL DOCUMENTS] below.
   Do not fabricate and do not infer legal content
   beyond the documents provided.
2. QUOTATION SCOPE: quote only the clause/point/passage
   that DIRECTLY answers the question. If an Article has
   many clauses but only 1-2 are relevant, quote exactly
   those clauses. Quoting a whole Article unnecessarily
   counts as a WRONG ANSWER (rambling, unfocused).
3. Cite the source as: According to [Document name],
   [Article/Clause]: "[verbatim relevant text]".
4. If the question can be answered directly by a short
   phrase (Yes/No, a number, a date, an authority name,
   a document name...), state that short answer FIRST,
   then quote the verbatim legal basis immediately after
   as evidence.
5. If [LEGAL DOCUMENTS] contains no provision relevant
   to the question, answer with exactly ONE sentence and
   nothing else: "No relevant provision was found in the
   documents provided."
\end{verbatim}

The user message concatenates the last three conversation turns (if any), the
formatted evidence block under the \texttt{[LEGAL DOCUMENTS]} header---with
internal article labels cleaned to the standard ``\textit{Article~N.\
<title>}'' form to avoid echoing internal identifiers into the
answer---followed by \texttt{[QUESTION]} and \texttt{[ANSWER]} markers.

\paragraph{Clarification prompt.}
When the final route is \texttt{clarify}, the clarification question is
generated under the following system prompt (translated into English), which
requires exactly one question and two to four concrete options:

\begin{verbatim}
You are a legal assistant. The user's question does not
yet specify a concrete legal document. Ask ONE short
clarifying question; do not ask more than once. MANDATORY:
you must suggest 2 to 4 concrete options for the user to
choose from. Each option must be on its own line, starting
with the exact format '- [Option] '.
\end{verbatim}

\paragraph{Stage~2 verifier contract.}
The Stage~2 verifier is prompted with the following system message and must
emit a single JSON object matching the schema; the \texttt{confidence}
attached to a \texttt{clarify} decision is the $p_c$ compared against
$\tau_{\mathrm{clar}}$ in Eq.~\eqref{eq:final_route}, and
\texttt{override\_stage1} records whether the verifier departed from the
Stage~1 proposal (the override rates of Table~\ref{tab:ambiguity}).

\begin{verbatim}
You are a Vietnamese legal-question analysis system.
Task: verify and refine the routing decision from
Stage 1.
You MUST return valid JSON in the schema below, with
no other text:
{
  "final_route":
    "<dense_retrieval|graph_traversal|
      hybrid_reasoning|clarify>",
  "confidence": <0.0-1.0>,
  "override_stage1": <true|false>,
  "override_reason": "<reason if override,
      null otherwise>",
  "complexity_level":
    "<simple|moderate|complex|ambiguous>",
  "reasoning_steps": [
    "<short analysis step 1>",
    "<short analysis step 2>"
  ],
  "sub_questions":
    ["<sub-question if multi-hop, [] otherwise>"],
  "ambiguity_flags": {
    "missing_entity": <true|false>,
    "pronoun_reference": <true|false>,
    "multi_interpretation": <true|false>,
    "incomplete_context": <true|false>
  },
  "clarify_question":
    "<clarifying question if final_route=clarify,
      null otherwise>",
  "resolved_referent":
    "<resolved referent if any, null otherwise>",
  "candidate_referents":
    [{"text": "<candidate>", "type": "<type>",
      "confidence": <0.0-1.0>}],
  "history_resolution_status":
    "<resolved|no_history|irrelevant_history|
      conflicting_history|not_needed>",
  "clarification_reason_type":
    "<missing_entity|multi_interpretation|
      incomplete_context|unresolved_history|
      conflicting_history|null>",
  "suggested_resolved_query":
    "<query with the referent replaced by the
      concrete entity if any, null otherwise>"
}
\end{verbatim}

\subsection{Training and Evaluation Configuration}\label{app:hyperparams}

Table~\ref{tab:hyperparams} lists the Stage~1 XGBoost configuration and the
baseline settings used in the cross-validated comparison of
Table~\ref{tab:routing_cv}; all classifiers share the same stratified folds
and balanced sample weights. Bootstrap confidence intervals use $B=1000$
percentile resamples with seed~$42$.

\begin{table}[!htbp]
\centering
\caption{Stage~1 training configuration and evaluation
protocol}\label{tab:hyperparams}
\footnotesize
\begin{tabular}{ll}
\toprule
\textbf{Setting} & \textbf{Value} \\
\midrule
\multicolumn{2}{l}{\emph{XGBoost (Stage~1, deployed)}}\\
\texttt{objective} & \texttt{multi:softprob} \\
\texttt{num\_class} & 4 (dense, graph, hybrid, clarify) \\
\texttt{n\_estimators} & 200 \\
\texttt{max\_depth} & 6 \\
\texttt{learning\_rate} & 0.05 \\
\texttt{subsample} & 0.8 \\
\texttt{colsample\_bytree} & 0.8 \\
\texttt{eval\_metric} & \texttt{mlogloss} \\
Sample weights & balanced per class \\
\midrule
\multicolumn{2}{l}{\emph{Baselines (identical folds)}}\\
Logistic Regression & $C=1.0$, \texttt{max\_iter}$=500$ \\
Random Forest & $200$ trees, balanced class weights \\
\midrule
\multicolumn{2}{l}{\emph{Protocol}}\\
Cross-validation & stratified 5-fold, shuffled, seed 42 \\
Training data & $800$ routing-labelled + $60$ synthetic clarify \\
Bootstrap CIs & $B=1000$, percentile, seed 42 \\
BERTScore encoder & \texttt{vinai/phobert-base}, \texttt{lang="vi"} \\
\bottomrule
\end{tabular}
\end{table}


% ---------------------------------------------------------------
\bibliography{biblio}

\end{document}
